\documentclass[10pt]{article}
\usepackage[a4paper,margin=20mm]{geometry}
\usepackage{graphicx,booktabs,longtable,amsmath,amssymb}
\usepackage[colorlinks=true,allcolors=blue]{hyperref}
\setlength{\emergencystretch}{1em}
\renewcommand{\arraystretch}{1.05}
\title{Supplementary Information\\[2mm]
\large A microbiome-informed Bayesian adaptation framework for personalized nutrient profiling}
\author{}\date{}
\begin{document}
\maketitle

All quantities below are participant-balanced unless stated otherwise, in units of
mmol\,l$^{-1}$\,h. Grouped cross-validation uses family membership where available and
participant identity otherwise, with a fixed seed. Confidence intervals are percentile intervals
from 2{,}000 paired participant-bootstrap resamples. Permutations number 200 and random
calibration draws number five. Sign tests exclude participants whose value is identical under both
models being compared.

\tableofcontents
\clearpage

\section*{Supplementary Note 1: Regression lock and two previously unreported facts}
\addcontentsline{toc}{section}{Note 1: Regression lock}

Eighteen quantities reported in the preceding version of this analysis were recomputed from the
saved out-of-fold predictions. All reproduce within a tolerance of $1.5\times10^{-3}$
(Supplementary Table~S1).

Two facts follow from the same recomputation and were not previously reported. First, the personal
response level accounts for 83.6\% of the total improvement over the shared predictor, so a
headline that quotes only the total conflates a level correction available to any method with the
composition-dependent adaptation that is specific to this representation. Second, applying the
conditional prior with no measured meals does not help. On the 902-meal query set the prior-only
prediction gives 2.3917 against 2.3639 for the shared predictor; on the 594-meal common query set
used for the calibration curve it gives 2.3610 against 2.3699. The direction depends on which
later meals are scored, so biological context provides no dependable benefit before the first
measured meal.

{\footnotesize\setlength{\tabcolsep}{3pt}
\begin{longtable}{lllll}
\toprule
\textbf{quantity} & \textbf{kind} & \textbf{published} & \textbf{recomputed} & \textbf{match} \\
\midrule
\endhead
shared & rmse & 2.3640 & 2.3639 & True \\
k1 & rmse & 2.1720 & 2.1720 & True \\
k3 & rmse & 1.9770 & 1.9772 & True \\
k5 & rmse & 1.9610 & 1.9612 & True \\
intercept\_only & rmse & 2.0410 & 2.0408 & True \\
nutrients\_only & rmse & 2.1640 & 2.1644 & True \\
intercept\_mac & rmse & 2.0100 & 2.0096 & True \\
intercept\_lipid & rmse & 2.0520 & 2.0521 & True \\
profile\_k3 & profile\_rmse & 1.8120 & 1.8123 & True \\
profile\_intercept\_only & profile\_rmse & 1.8810 & 1.8813 & True \\
profile\_nutrients\_only & profile\_rmse & 1.8530 & 1.8532 & True \\
profile\_intercept\_mac & profile\_rmse & 1.8220 & 1.8220 & True \\
profile\_intercept\_lipid & profile\_rmse & 1.8850 & 1.8852 & True \\
profile\_shared & profile\_rmse & 1.8810 & 1.8813 & True \\
micro\_vs\_population\_k3 & difference & 0.0061 & 0.0061 & True \\
micro\_vs\_host\_k3 & difference & 0.0042 & 0.0042 & True \\
relative\_shared\_to\_full & pct & 16.4000 & 16.3560 & True \\
relative\_intercept\_to\_full & pct & 3.1000 & 3.1140 & True \\
\bottomrule
\end{longtable}}

\clearpage

\section*{Supplementary Note 2: Three-layer decomposition and the primary estimand}
\addcontentsline{toc}{section}{Note 2: Three-layer decomposition}

The primary estimand is the reduction in participant-balanced RMSE of the full three-coefficient
model relative to intercept-only calibration at three calibration meals. The reduction relative to
the shared predictor is reported for context only.

In the free-living cohort the shared predictor reaches 2.364, intercept-only calibration 2.041 and
the full model 1.977 (Supplementary Table~S2). The channel increment is 0.064
(95\% CI, 0.022--0.107), a relative reduction of 3.11\% (1.15--4.89), with lower error in 33 of 44
participants (sign test $P=0.0013$). An unshrunk personal mean reaches 2.121, so the
empirical-Bayes prior contributes 0.080 (95\% CI, 0.016--0.154) on its own. The lipid channel is
not supported in this cohort (0.032; 95\% CI, $-0.007$ to 0.079; $P=0.45$), whereas the
carbohydrate channel is (0.075; 95\% CI, 0.037--0.117; $P=1.1\times10^{-4}$). Supplementary
Table~S3 gives every contrast on both metrics.

The bootstrap standard error of the primary estimand is 0.0216, giving a minimum detectable effect
at 80\% power of 0.0605 against an observed effect of 0.0636. This cohort is therefore only just
powered for its own primary comparison and has no power at all for effects of order 0.006.

{\footnotesize\setlength{\tabcolsep}{3pt}
\begin{longtable}{lllll}
\toprule
\textbf{model} & \textbf{meals} & \textbf{people} & \textbf{rmse} & \textbf{profile rmse} \\
\midrule
\endhead
full & 902 & 44 & 1.9772 & 1.8123 \\
intercept\_plus\_mac & 902 & 44 & 2.0096 & 1.8220 \\
intercept\_only & 902 & 44 & 2.0408 & 1.8813 \\
intercept\_plus\_lipid & 902 & 44 & 2.0521 & 1.8852 \\
raw\_personal\_mean\_k3 & 902 & 44 & 2.1211 & 1.8813 \\
channels\_without\_intercept & 902 & 44 & 2.1644 & 1.8532 \\
shared & 902 & 44 & 2.3639 & 1.8813 \\
\bottomrule
\end{longtable}}

{\footnotesize\setlength{\tabcolsep}{3pt}
\begin{longtable}{llllllllll}
\toprule
\textbf{contrast} & \textbf{metric} & \textbf{method} & \textbf{reference} & \textbf{difference} & \textbf{diff lo} & \textbf{diff hi} & \textbf{relative pct} & \textbf{better} & \textbf{sign p} \\
\midrule
\endhead
level\_correction & rmse & 2.0408 & 2.3639 & 0.3231 & 0.1562 & 0.4887 & 13.6677 & 30 & 0.0226 \\
level\_correction & profile\_rmse & 1.8813 & 1.8813 & 0.0000 & -0.0000 & 0.0000 & 0.0000 & 14 & 0.0226 \\
channel\_increment\_PRIMARY & rmse & 1.9772 & 2.0408 & 0.0636 & 0.0225 & 0.1066 & 3.1142 & 33 & 0.0013 \\
channel\_increment\_PRIMARY & profile\_rmse & 1.8123 & 1.8813 & 0.0690 & 0.0238 & 0.1192 & 3.6695 & 29 & 0.0488 \\
total\_gain\_context & rmse & 1.9772 & 2.3639 & 0.3866 & 0.2042 & 0.5639 & 16.3563 & 30 & 0.0226 \\
total\_gain\_context & profile\_rmse & 1.8123 & 1.8813 & 0.0690 & 0.0238 & 0.1192 & 3.6695 & 29 & 0.0488 \\
shrinkage\_vs\_raw\_mean & rmse & 2.0408 & 2.1211 & 0.0804 & 0.0160 & 0.1536 & 3.7881 & 30 & 0.0226 \\
shrinkage\_vs\_raw\_mean & profile\_rmse & 1.8813 & 1.8813 & 0.0000 & -0.0000 & 0.0000 & 0.0000 & 11 & 0.0013 \\
full\_vs\_raw\_mean & rmse & 1.9772 & 2.1211 & 0.1439 & 0.0717 & 0.2208 & 6.7843 & 36 & 0.0000 \\
full\_vs\_raw\_mean & profile\_rmse & 1.8123 & 1.8813 & 0.0690 & 0.0238 & 0.1192 & 3.6695 & 29 & 0.0488 \\
mac\_contribution & rmse & 1.9772 & 2.0521 & 0.0748 & 0.0369 & 0.1165 & 3.6470 & 35 & 0.0001 \\
mac\_contribution & profile\_rmse & 1.8123 & 1.8852 & 0.0729 & 0.0273 & 0.1229 & 3.8667 & 29 & 0.0488 \\
lipid\_contribution & rmse & 1.9772 & 2.0096 & 0.0324 & -0.0072 & 0.0787 & 1.6105 & 25 & 0.4514 \\
lipid\_contribution & profile\_rmse & 1.8123 & 1.8220 & 0.0098 & -0.0049 & 0.0254 & 0.5352 & 28 & 0.0961 \\
intercept\_necessity & rmse & 1.9772 & 2.1644 & 0.1872 & 0.0706 & 0.3055 & 8.6471 & 27 & 0.1742 \\
intercept\_necessity & profile\_rmse & 1.8123 & 1.8532 & 0.0409 & 0.0104 & 0.0757 & 2.2072 & 28 & 0.0961 \\
\bottomrule
\end{longtable}}

\clearpage

\section*{Supplementary Note 3: Prior arms, nested selection and expert review}
\addcontentsline{toc}{section}{Note 3: Prior arms and selection}

Conditional priors are grouped into four arms that do not overlap: population, host, microbial and
an exploratory host-plus-microbial arm. Selection runs separately inside each arm and the chosen
configuration is asserted to remain in its arm. Without this discipline a microbial arm can be
occupied by a host-plus-microbial representation, which makes any microbial increment
uninterpretable. Supplementary Table~S7 lists the configuration selected in every outer fold and
arm; the microbial arm selected a taxon representation in four folds and a functional
representation in one.

Six nutrition experts reviewed the nutrient-channel assignments and their supporting evidence in
two rounds, before the adaptive analysis was run. In the second round 176 of 180 assignment
judgements were rated reasonable and 169 of 180 evidence explanations were rated sufficient. A
separate review of food-group explanations rated 68 of 72 reasonable. A third review of case
communication rated 64 of 96 reasonable, that is 66.7\% (95\% CI, 56.8--75.3\%), which is a
negative finding and is reported as one.

{\footnotesize\setlength{\tabcolsep}{3pt}
\begin{longtable}{llll}
\toprule
\textbf{fold} & \textbf{arm} & \textbf{config} & \textbf{inner rmse} \\
\midrule
\endhead
0 & population & masked\_uniform\_noise3.0 & 0.3001 \\
0 & host & masked\_host\_bw1.0\_mean0.5 & 0.3008 \\
0 & micro & masked\_taxa\_bw0.3\_mean0.5 & 0.3006 \\
0 & host\_micro & masked\_host\_taxa\_bw0.3\_mean0.0 & 0.3008 \\
1 & population & masked\_uniform\_noise3.0 & 0.3027 \\
1 & host & masked\_host\_bw0.3\_mean0.5 & 0.3007 \\
1 & micro & masked\_taxa\_bw1.0\_mean0.5 & 0.3008 \\
1 & host\_micro & masked\_host\_function\_bw0.3\_mean0.5 & 0.3007 \\
2 & population & masked\_uniform\_noise3.0 & 0.2916 \\
2 & host & masked\_host\_bw1.0\_mean0.0 & 0.2920 \\
2 & micro & masked\_taxa\_bw0.3\_mean0.0 & 0.2920 \\
2 & host\_micro & masked\_host\_taxa\_bw1.0\_mean0.0 & 0.2920 \\
3 & population & masked\_uniform\_noise3.0 & 0.3044 \\
3 & host & masked\_host\_bw1.0\_mean0.0 & 0.3040 \\
3 & micro & masked\_taxa\_bw0.3\_mean0.5 & 0.3038 \\
3 & host\_micro & masked\_host\_taxa\_bw0.3\_mean0.5 & 0.3039 \\
4 & population & masked\_uniform\_noise3.0 & 0.3008 \\
4 & host & masked\_host\_bw3.0\_mean0.5 & 0.3007 \\
4 & micro & masked\_function\_bw3.0\_mean0.0 & 0.3007 \\
4 & host\_micro & masked\_host\_function\_bw1.0\_mean0.5 & 0.3007 \\
\bottomrule
\end{longtable}}

\clearpage

\section*{Supplementary Note 4: Negative control for microbial conditioning}
\addcontentsline{toc}{section}{Note 4: Negative control}

The microbiome enters the model in two roles: it shapes the prior estimated from training
participants, and it selects which training participants a new participant resembles. The negative
control breaks both. Every participant, in training and in held-out data, receives another
participant's complete vector of 2{,}001 microbial and functional features. The empirical-Bayes
moments and the conditional prior are then re-estimated and the evaluation repeated. Two hundred
permutations were drawn.

The observed advantage of microbial conditioning over an unconditional prior is $+0.00216$. The
permutation reference distribution has median $+0.00229$ and a 2.5th to 97.5th percentile range of
0.00189 to 0.00269, so the observed value sits at the 27.5th percentile with a permutation
$P$ of 0.73. The observed microbial-arm RMSE of 1.9812 is matched by a permuted median of 1.9811.
Breaking the true microbiome-to-participant correspondence therefore changes nothing. What the
conditional prior gains is attributable to weighting training participants at all, not to the
biological variable used for the weighting.

{\footnotesize\setlength{\tabcolsep}{3pt}
\begin{longtable}{llllllll}
\toprule
\textbf{metric} & \textbf{observed} & \textbf{perm median} & \textbf{perm q025} & \textbf{perm q975} & \textbf{perm quantile of observed} & \textbf{n perm at or above} & \textbf{permutation p} \\
\midrule
\endhead
micro\_vs\_population & 0.00216 & 0.00229 & 0.00189 & 0.00269 & 0.27500 & 145 & 0.72637 \\
micro\_vs\_host & -0.00183 & -0.00170 & -0.00211 & -0.00130 & 0.27500 & 145 & 0.72637 \\
\bottomrule
\end{longtable}}

\clearpage

\section*{Supplementary Note 5: External replication and frozen transport}
\addcontentsline{toc}{section}{Note 5: External replication and transport}

The procedure was repeated in full, including nested selection, in 998 participants eating seven
standardized meal compositions. Calibration meals were drawn at random because these meals carry
no time order, and five independent draws were evaluated.

Beyond the personal offset the channels reduce error by 4.54\% (95\% CI, 3.71--5.39), with lower
error in 615 of 998 participants. Both channels are supported, the lipid channel by 3.44\%
(95\% CI, 0.0083--0.0131) and the carbohydrate channel by 1.21\% (0.0029--0.0045). Against an
unshrunk personal mean the full model gains 7.29\%.

For transport, the shared predictor and the prior were fitted on the 998-participant cohort,
frozen, and applied to a geographically separate cohort of 100 participants. No target-cohort data
entered training or selection, and disjointness was asserted. The channel increment is 6.92\%
(95\% CI, 0.0129--0.0327; $P=0.005$). The personal response level does not transport
(0.0010; 95\% CI, $-0.019$ to 0.019; $P=0.46$). Supplementary Tables~S4 and~S5 give all models and
contrasts for both cohorts.

{\footnotesize\setlength{\tabcolsep}{3pt}
\begin{longtable}{llllllll}
\toprule
\textbf{cohort} & \textbf{arm} & \textbf{model} & \textbf{meals} & \textbf{people} & \textbf{rmse} & \textbf{rmse sd across draws} & \textbf{profile rmse} \\
\midrule
\endhead
PREDICT\_UK & host & channels\_only & 3987.0000 & 998.0000 & 0.3342 & 0.0032 & 0.2674 \\
PREDICT\_UK & host & full & 3987.0000 & 998.0000 & 0.3003 & 0.0033 & 0.2396 \\
PREDICT\_UK & host & intercept\_only & 3987.0000 & 998.0000 & 0.3148 & 0.0031 & 0.2462 \\
PREDICT\_UK & host & intercept\_plus\_lipid & 3987.0000 & 998.0000 & 0.3040 & 0.0036 & 0.2417 \\
PREDICT\_UK & host & intercept\_plus\_mac & 3987.0000 & 998.0000 & 0.3112 & 0.0028 & 0.2444 \\
PREDICT\_UK & host & raw\_personal\_mean & 3987.0000 & 998.0000 & 0.3241 & 0.0033 & 0.2462 \\
PREDICT\_UK & host & shared & 3987.0000 & 998.0000 & 0.3298 & 0.0034 & 0.2462 \\
PREDICT\_UK & host\_micro & channels\_only & 3987.0000 & 998.0000 & 0.3340 & 0.0032 & 0.2673 \\
PREDICT\_UK & host\_micro & full & 3987.0000 & 998.0000 & 0.3003 & 0.0034 & 0.2396 \\
PREDICT\_UK & host\_micro & intercept\_only & 3987.0000 & 998.0000 & 0.3148 & 0.0031 & 0.2462 \\
PREDICT\_UK & host\_micro & intercept\_plus\_lipid & 3987.0000 & 998.0000 & 0.3041 & 0.0036 & 0.2417 \\
PREDICT\_UK & host\_micro & intercept\_plus\_mac & 3987.0000 & 998.0000 & 0.3112 & 0.0028 & 0.2444 \\
PREDICT\_UK & host\_micro & raw\_personal\_mean & 3987.0000 & 998.0000 & 0.3241 & 0.0033 & 0.2462 \\
PREDICT\_UK & host\_micro & shared & 3987.0000 & 998.0000 & 0.3298 & 0.0034 & 0.2462 \\
PREDICT\_UK & micro & channels\_only & 3987.0000 & 998.0000 & 0.3340 & 0.0032 & 0.2673 \\
PREDICT\_UK & micro & full & 3987.0000 & 998.0000 & 0.3004 & 0.0034 & 0.2396 \\
PREDICT\_UK & micro & intercept\_only & 3987.0000 & 998.0000 & 0.3147 & 0.0031 & 0.2462 \\
PREDICT\_UK & micro & intercept\_plus\_lipid & 3987.0000 & 998.0000 & 0.3041 & 0.0037 & 0.2417 \\
PREDICT\_UK & micro & intercept\_plus\_mac & 3987.0000 & 998.0000 & 0.3111 & 0.0028 & 0.2444 \\
PREDICT\_UK & micro & raw\_personal\_mean & 3987.0000 & 998.0000 & 0.3241 & 0.0033 & 0.2462 \\
PREDICT\_UK & micro & shared & 3987.0000 & 998.0000 & 0.3298 & 0.0034 & 0.2462 \\
PREDICT\_UK & population & channels\_only & 3987.0000 & 998.0000 & 0.3260 & 0.0029 & 0.2579 \\
PREDICT\_UK & population & full & 3987.0000 & 998.0000 & 0.3023 & 0.0035 & 0.2422 \\
PREDICT\_UK & population & intercept\_only & 3987.0000 & 998.0000 & 0.3076 & 0.0032 & 0.2462 \\
PREDICT\_UK & population & intercept\_plus\_lipid & 3987.0000 & 998.0000 & 0.3029 & 0.0036 & 0.2431 \\
PREDICT\_UK & population & intercept\_plus\_mac & 3987.0000 & 998.0000 & 0.3050 & 0.0031 & 0.2448 \\
PREDICT\_UK & population & raw\_personal\_mean & 3987.0000 & 998.0000 & 0.3241 & 0.0033 & 0.2462 \\
PREDICT\_UK & population & shared & 3987.0000 & 998.0000 & 0.3298 & 0.0034 & 0.2462 \\
PREDICT\_US\_transport & host & channels\_only & 399.0000 & 100.0000 & 0.3343 & 0.0091 & 0.2645 \\
PREDICT\_US\_transport & host & full & 399.0000 & 100.0000 & 0.3000 & 0.0052 & 0.2356 \\
PREDICT\_US\_transport & host & intercept\_only & 399.0000 & 100.0000 & 0.3227 & 0.0050 & 0.2438 \\
PREDICT\_US\_transport & host & intercept\_plus\_lipid & 399.0000 & 100.0000 & 0.3051 & 0.0054 & 0.2380 \\
PREDICT\_US\_transport & host & intercept\_plus\_mac & 399.0000 & 100.0000 & 0.3184 & 0.0047 & 0.2424 \\
PREDICT\_US\_transport & host & raw\_personal\_mean & 399.0000 & 100.0000 & 0.3357 & 0.0059 & 0.2438 \\
PREDICT\_US\_transport & host & shared & 399.0000 & 100.0000 & 0.3231 & 0.0019 & 0.2438 \\
PREDICT\_US\_transport & host\_micro & channels\_only & 399.0000 & 100.0000 & 0.3341 & 0.0090 & 0.2643 \\
PREDICT\_US\_transport & host\_micro & full & 399.0000 & 100.0000 & 0.3001 & 0.0052 & 0.2359 \\
PREDICT\_US\_transport & host\_micro & intercept\_only & 399.0000 & 100.0000 & 0.3223 & 0.0050 & 0.2438 \\
PREDICT\_US\_transport & host\_micro & intercept\_plus\_lipid & 399.0000 & 100.0000 & 0.3051 & 0.0053 & 0.2382 \\
PREDICT\_US\_transport & host\_micro & intercept\_plus\_mac & 399.0000 & 100.0000 & 0.3180 & 0.0047 & 0.2424 \\
PREDICT\_US\_transport & host\_micro & raw\_personal\_mean & 399.0000 & 100.0000 & 0.3357 & 0.0059 & 0.2438 \\
PREDICT\_US\_transport & host\_micro & shared & 399.0000 & 100.0000 & 0.3231 & 0.0019 & 0.2438 \\
PREDICT\_US\_transport & micro & channels\_only & 399.0000 & 100.0000 & 0.3336 & 0.0088 & 0.2639 \\
PREDICT\_US\_transport & micro & full & 399.0000 & 100.0000 & 0.2998 & 0.0050 & 0.2357 \\
PREDICT\_US\_transport & micro & intercept\_only & 399.0000 & 100.0000 & 0.3221 & 0.0050 & 0.2438 \\
PREDICT\_US\_transport & micro & intercept\_plus\_lipid & 399.0000 & 100.0000 & 0.3048 & 0.0053 & 0.2381 \\
PREDICT\_US\_transport & micro & intercept\_plus\_mac & 399.0000 & 100.0000 & 0.3179 & 0.0047 & 0.2424 \\
PREDICT\_US\_transport & micro & raw\_personal\_mean & 399.0000 & 100.0000 & 0.3357 & 0.0059 & 0.2438 \\
PREDICT\_US\_transport & micro & shared & 399.0000 & 100.0000 & 0.3231 & 0.0019 & 0.2438 \\
PREDICT\_US\_transport & population & channels\_only & 399.0000 & 100.0000 & 0.3228 & 0.0052 & 0.2551 \\
PREDICT\_US\_transport & population & full & 399.0000 & 100.0000 & 0.2997 & 0.0039 & 0.2389 \\
PREDICT\_US\_transport & population & intercept\_only & 399.0000 & 100.0000 & 0.3106 & 0.0040 & 0.2438 \\
PREDICT\_US\_transport & population & intercept\_plus\_lipid & 399.0000 & 100.0000 & 0.3013 & 0.0039 & 0.2400 \\
PREDICT\_US\_transport & population & intercept\_plus\_mac & 399.0000 & 100.0000 & 0.3073 & 0.0038 & 0.2426 \\
PREDICT\_US\_transport & population & raw\_personal\_mean & 399.0000 & 100.0000 & 0.3357 & 0.0059 & 0.2438 \\
PREDICT\_US\_transport & population & shared & 399.0000 & 100.0000 & 0.3231 & 0.0019 & 0.2438 \\
\bottomrule
\end{longtable}}

{\footnotesize\setlength{\tabcolsep}{3pt}
\begin{longtable}{llllllllll}
\toprule
\textbf{cohort} & \textbf{arm} & \textbf{contrast} & \textbf{metric} & \textbf{difference} & \textbf{diff lo} & \textbf{diff hi} & \textbf{relative pct} & \textbf{better} & \textbf{sign p} \\
\midrule
\endhead
PREDICT\_UK & host & channel\_increment\_PRIMARY & profile\_rmse & 0.0066 & 0.0049 & 0.0083 & 2.6859 & 577.8000 & 0.0000 \\
PREDICT\_UK & host & channel\_increment\_PRIMARY & rmse & 0.0146 & 0.0119 & 0.0172 & 4.6265 & 615.8000 & 0.0000 \\
PREDICT\_UK & host & full\_vs\_raw\_mean & profile\_rmse & 0.0066 & 0.0049 & 0.0083 & 2.6859 & 577.8000 & 0.0000 \\
PREDICT\_UK & host & full\_vs\_raw\_mean & rmse & 0.0238 & 0.0198 & 0.0279 & 7.3481 & 617.0000 & 0.0000 \\
PREDICT\_UK & host & level\_correction & profile\_rmse & 0.0000 & -0.0000 & 0.0000 & 0.0000 & 394.0000 & 0.0000 \\
PREDICT\_UK & host & level\_correction & rmse & 0.0150 & 0.0095 & 0.0205 & 4.5487 & 596.4000 & 0.0000 \\
PREDICT\_UK & host & lipid\_contribution & profile\_rmse & 0.0049 & 0.0033 & 0.0065 & 1.9869 & 558.2000 & 0.0006 \\
PREDICT\_UK & host & lipid\_contribution & rmse & 0.0109 & 0.0085 & 0.0133 & 3.5112 & 592.2000 & 0.0000 \\
PREDICT\_UK & host & mac\_contribution & profile\_rmse & 0.0021 & 0.0016 & 0.0026 & 0.8744 & 572.4000 & 0.0000 \\
PREDICT\_UK & host & mac\_contribution & rmse & 0.0038 & 0.0030 & 0.0046 & 1.2415 & 579.2000 & 0.0000 \\
PREDICT\_UK & host & shrinkage\_vs\_raw\_mean & profile\_rmse & -0.0000 & -0.0000 & 0.0000 & -0.0000 & 396.4000 & 0.0000 \\
PREDICT\_UK & host & shrinkage\_vs\_raw\_mean & rmse & 0.0092 & 0.0076 & 0.0109 & 2.8546 & 576.4000 & 0.0000 \\
PREDICT\_UK & host & total\_gain\_context & profile\_rmse & 0.0066 & 0.0049 & 0.0083 & 2.6859 & 577.8000 & 0.0000 \\
PREDICT\_UK & host & total\_gain\_context & rmse & 0.0296 & 0.0249 & 0.0342 & 8.9671 & 648.4000 & 0.0000 \\
PREDICT\_UK & host\_micro & channel\_increment\_PRIMARY & profile\_rmse & 0.0066 & 0.0049 & 0.0083 & 2.6691 & 577.2000 & 0.0000 \\
PREDICT\_UK & host\_micro & channel\_increment\_PRIMARY & rmse & 0.0145 & 0.0118 & 0.0171 & 4.5961 & 612.4000 & 0.0000 \\
PREDICT\_UK & host\_micro & full\_vs\_raw\_mean & profile\_rmse & 0.0066 & 0.0049 & 0.0083 & 2.6691 & 577.2000 & 0.0000 \\
PREDICT\_UK & host\_micro & full\_vs\_raw\_mean & rmse & 0.0238 & 0.0198 & 0.0278 & 7.3273 & 613.6000 & 0.0000 \\
PREDICT\_UK & host\_micro & level\_correction & profile\_rmse & 0.0000 & -0.0000 & 0.0000 & 0.0000 & 384.0000 & 0.0000 \\
PREDICT\_UK & host\_micro & level\_correction & rmse & 0.0150 & 0.0095 & 0.0205 & 4.5579 & 595.4000 & 0.0000 \\
PREDICT\_UK & host\_micro & lipid\_contribution & profile\_rmse & 0.0048 & 0.0032 & 0.0064 & 1.9694 & 558.0000 & 0.0005 \\
PREDICT\_UK & host\_micro & lipid\_contribution & rmse & 0.0108 & 0.0084 & 0.0132 & 3.4829 & 589.8000 & 0.0000 \\
PREDICT\_UK & host\_micro & mac\_contribution & profile\_rmse & 0.0021 & 0.0016 & 0.0026 & 0.8706 & 572.0000 & 0.0000 \\
PREDICT\_UK & host\_micro & mac\_contribution & rmse & 0.0037 & 0.0029 & 0.0046 & 1.2322 & 578.2000 & 0.0000 \\
PREDICT\_UK & host\_micro & shrinkage\_vs\_raw\_mean & profile\_rmse & 0.0000 & -0.0000 & 0.0000 & 0.0000 & 386.2000 & 0.0000 \\
PREDICT\_UK & host\_micro & shrinkage\_vs\_raw\_mean & rmse & 0.0093 & 0.0077 & 0.0109 & 2.8638 & 578.8000 & 0.0000 \\
PREDICT\_UK & host\_micro & total\_gain\_context & profile\_rmse & 0.0066 & 0.0049 & 0.0083 & 2.6691 & 577.2000 & 0.0000 \\
PREDICT\_UK & host\_micro & total\_gain\_context & rmse & 0.0295 & 0.0248 & 0.0341 & 8.9468 & 648.4000 & 0.0000 \\
PREDICT\_UK & micro & channel\_increment\_PRIMARY & profile\_rmse & 0.0065 & 0.0049 & 0.0083 & 2.6584 & 574.8000 & 0.0000 \\
PREDICT\_UK & micro & channel\_increment\_PRIMARY & rmse & 0.0143 & 0.0117 & 0.0169 & 4.5414 & 614.6000 & 0.0000 \\
PREDICT\_UK & micro & full\_vs\_raw\_mean & profile\_rmse & 0.0065 & 0.0049 & 0.0083 & 2.6584 & 574.8000 & 0.0000 \\
PREDICT\_UK & micro & full\_vs\_raw\_mean & rmse & 0.0236 & 0.0197 & 0.0277 & 7.2890 & 614.0000 & 0.0000 \\
PREDICT\_UK & micro & level\_correction & profile\_rmse & 0.0000 & -0.0000 & 0.0000 & 0.0000 & 378.2000 & 0.0000 \\
PREDICT\_UK & micro & level\_correction & rmse & 0.0151 & 0.0096 & 0.0206 & 4.5732 & 595.4000 & 0.0000 \\
PREDICT\_UK & micro & lipid\_contribution & profile\_rmse & 0.0048 & 0.0032 & 0.0064 & 1.9648 & 559.6000 & 0.0009 \\
PREDICT\_UK & micro & lipid\_contribution & rmse & 0.0107 & 0.0083 & 0.0131 & 3.4362 & 589.4000 & 0.0000 \\
PREDICT\_UK & micro & mac\_contribution & profile\_rmse & 0.0021 & 0.0016 & 0.0026 & 0.8604 & 575.6000 & 0.0000 \\
PREDICT\_UK & micro & mac\_contribution & rmse & 0.0037 & 0.0029 & 0.0045 & 1.2077 & 581.4000 & 0.0000 \\
PREDICT\_UK & micro & shrinkage\_vs\_raw\_mean & profile\_rmse & -0.0000 & -0.0000 & 0.0000 & -0.0000 & 390.2000 & 0.0000 \\
PREDICT\_UK & micro & shrinkage\_vs\_raw\_mean & rmse & 0.0093 & 0.0077 & 0.0110 & 2.8793 & 578.8000 & 0.0000 \\
PREDICT\_UK & micro & total\_gain\_context & profile\_rmse & 0.0065 & 0.0049 & 0.0083 & 2.6584 & 574.8000 & 0.0000 \\
PREDICT\_UK & micro & total\_gain\_context & rmse & 0.0294 & 0.0247 & 0.0340 & 8.9092 & 649.6000 & 0.0000 \\
PREDICT\_UK & population & channel\_increment\_PRIMARY & profile\_rmse & 0.0039 & 0.0031 & 0.0047 & 1.6007 & 594.0000 & 0.0000 \\
PREDICT\_UK & population & channel\_increment\_PRIMARY & rmse & 0.0054 & 0.0037 & 0.0071 & 1.7443 & 546.8000 & 0.0353 \\
PREDICT\_UK & population & full\_vs\_raw\_mean & profile\_rmse & 0.0039 & 0.0031 & 0.0047 & 1.6007 & 594.0000 & 0.0000 \\
PREDICT\_UK & population & full\_vs\_raw\_mean & rmse & 0.0218 & 0.0169 & 0.0269 & 6.7295 & 543.8000 & 0.0229 \\
PREDICT\_UK & population & level\_correction & profile\_rmse & 0.0000 & -0.0000 & 0.0000 & 0.0000 & 381.4000 & 0.0000 \\
PREDICT\_UK & population & level\_correction & rmse & 0.0222 & 0.0181 & 0.0262 & 6.7340 & 641.0000 & 0.0000 \\
PREDICT\_UK & population & lipid\_contribution & profile\_rmse & 0.0026 & 0.0018 & 0.0033 & 1.0549 & 563.8000 & 0.0004 \\
PREDICT\_UK & population & lipid\_contribution & rmse & 0.0028 & 0.0014 & 0.0042 & 0.9160 & 524.2000 & 0.1749 \\
PREDICT\_UK & population & mac\_contribution & profile\_rmse & 0.0009 & 0.0006 & 0.0012 & 0.3701 & 552.6000 & 0.0018 \\
PREDICT\_UK & population & mac\_contribution & rmse & 0.0006 & 0.0003 & 0.0010 & 0.2076 & 495.4000 & 0.7136 \\
PREDICT\_UK & population & shrinkage\_vs\_raw\_mean & profile\_rmse & -0.0000 & -0.0000 & 0.0000 & -0.0000 & 395.6000 & 0.0000 \\
PREDICT\_UK & population & shrinkage\_vs\_raw\_mean & rmse & 0.0164 & 0.0130 & 0.0200 & 5.0753 & 547.2000 & 0.0189 \\
PREDICT\_UK & population & total\_gain\_context & profile\_rmse & 0.0039 & 0.0031 & 0.0047 & 1.6007 & 594.0000 & 0.0000 \\
PREDICT\_UK & population & total\_gain\_context & rmse & 0.0276 & 0.0245 & 0.0306 & 8.3620 & 695.4000 & 0.0000 \\
PREDICT\_US\_transport & host & channel\_increment\_PRIMARY & profile\_rmse & 0.0082 & 0.0027 & 0.0137 & 3.3586 & 57.8000 & 0.2375 \\
PREDICT\_US\_transport & host & channel\_increment\_PRIMARY & rmse & 0.0227 & 0.0131 & 0.0335 & 7.0215 & 66.2000 & 0.0026 \\
PREDICT\_US\_transport & host & full\_vs\_raw\_mean & profile\_rmse & 0.0082 & 0.0027 & 0.0137 & 3.3586 & 57.8000 & 0.2375 \\
PREDICT\_US\_transport & host & full\_vs\_raw\_mean & rmse & 0.0357 & 0.0211 & 0.0519 & 10.6236 & 66.2000 & 0.0123 \\
PREDICT\_US\_transport & host & level\_correction & profile\_rmse & 0.0000 & -0.0000 & 0.0000 & 0.0000 & 38.6000 & 0.1042 \\
PREDICT\_US\_transport & host & level\_correction & rmse & 0.0005 & -0.0197 & 0.0191 & 0.1440 & 54.6000 & 0.4547 \\
PREDICT\_US\_transport & host & lipid\_contribution & profile\_rmse & 0.0068 & 0.0017 & 0.0119 & 2.7925 & 60.2000 & 0.1333 \\
PREDICT\_US\_transport & host & lipid\_contribution & rmse & 0.0184 & 0.0100 & 0.0277 & 5.7887 & 65.2000 & 0.0127 \\
PREDICT\_US\_transport & host & mac\_contribution & profile\_rmse & 0.0024 & 0.0003 & 0.0045 & 1.0128 & 57.4000 & 0.2527 \\
PREDICT\_US\_transport & host & mac\_contribution & rmse & 0.0051 & 0.0020 & 0.0085 & 1.6665 & 56.6000 & 0.2704 \\
PREDICT\_US\_transport & host & shrinkage\_vs\_raw\_mean & profile\_rmse & 0.0000 & -0.0000 & 0.0000 & -0.0000 & 36.8000 & 0.0268 \\
PREDICT\_US\_transport & host & shrinkage\_vs\_raw\_mean & rmse & 0.0130 & 0.0074 & 0.0189 & 3.8753 & 61.8000 & 0.0889 \\
PREDICT\_US\_transport & host & total\_gain\_context & profile\_rmse & 0.0082 & 0.0027 & 0.0137 & 3.3586 & 57.8000 & 0.2375 \\
PREDICT\_US\_transport & host & total\_gain\_context & rmse & 0.0231 & 0.0081 & 0.0378 & 7.1587 & 62.2000 & 0.0361 \\
PREDICT\_US\_transport & host\_micro & channel\_increment\_PRIMARY & profile\_rmse & 0.0079 & 0.0025 & 0.0134 & 3.2415 & 58.0000 & 0.2332 \\
PREDICT\_US\_transport & host\_micro & channel\_increment\_PRIMARY & rmse & 0.0222 & 0.0127 & 0.0329 & 6.8839 & 65.6000 & 0.0032 \\
PREDICT\_US\_transport & host\_micro & full\_vs\_raw\_mean & profile\_rmse & 0.0079 & 0.0025 & 0.0134 & 3.2415 & 58.0000 & 0.2332 \\
PREDICT\_US\_transport & host\_micro & full\_vs\_raw\_mean & rmse & 0.0356 & 0.0210 & 0.0520 & 10.5999 & 65.6000 & 0.0130 \\
PREDICT\_US\_transport & host\_micro & level\_correction & profile\_rmse & 0.0000 & -0.0000 & 0.0000 & -0.0000 & 40.2000 & 0.2419 \\
PREDICT\_US\_transport & host\_micro & level\_correction & rmse & 0.0009 & -0.0192 & 0.0193 & 0.2648 & 54.6000 & 0.4547 \\
PREDICT\_US\_transport & host\_micro & lipid\_contribution & profile\_rmse & 0.0065 & 0.0015 & 0.0116 & 2.6783 & 60.2000 & 0.1178 \\
PREDICT\_US\_transport & host\_micro & lipid\_contribution & rmse & 0.0180 & 0.0097 & 0.0271 & 5.6447 & 64.6000 & 0.0198 \\
PREDICT\_US\_transport & host\_micro & mac\_contribution & profile\_rmse & 0.0023 & 0.0003 & 0.0045 & 0.9818 & 57.0000 & 0.2913 \\
PREDICT\_US\_transport & host\_micro & mac\_contribution & rmse & 0.0050 & 0.0019 & 0.0084 & 1.6295 & 56.6000 & 0.2926 \\
PREDICT\_US\_transport & host\_micro & shrinkage\_vs\_raw\_mean & profile\_rmse & -0.0000 & -0.0000 & 0.0000 & -0.0000 & 40.6000 & 0.1081 \\
PREDICT\_US\_transport & host\_micro & shrinkage\_vs\_raw\_mean & rmse & 0.0134 & 0.0076 & 0.0195 & 3.9916 & 61.8000 & 0.0889 \\
PREDICT\_US\_transport & host\_micro & total\_gain\_context & profile\_rmse & 0.0079 & 0.0025 & 0.0134 & 3.2415 & 58.0000 & 0.2332 \\
PREDICT\_US\_transport & host\_micro & total\_gain\_context & rmse & 0.0231 & 0.0083 & 0.0374 & 7.1338 & 62.2000 & 0.0361 \\
PREDICT\_US\_transport & micro & channel\_increment\_PRIMARY & profile\_rmse & 0.0080 & 0.0026 & 0.0135 & 3.3049 & 58.8000 & 0.1714 \\
PREDICT\_US\_transport & micro & channel\_increment\_PRIMARY & rmse & 0.0223 & 0.0129 & 0.0327 & 6.9198 & 65.6000 & 0.0053 \\
PREDICT\_US\_transport & micro & full\_vs\_raw\_mean & profile\_rmse & 0.0080 & 0.0026 & 0.0135 & 3.3049 & 58.8000 & 0.1714 \\
PREDICT\_US\_transport & micro & full\_vs\_raw\_mean & rmse & 0.0359 & 0.0213 & 0.0520 & 10.6799 & 67.0000 & 0.0079 \\
PREDICT\_US\_transport & micro & level\_correction & profile\_rmse & -0.0000 & -0.0000 & 0.0000 & -0.0000 & 40.4000 & 0.1385 \\
PREDICT\_US\_transport & micro & level\_correction & rmse & 0.0010 & -0.0190 & 0.0194 & 0.3164 & 54.8000 & 0.4163 \\
PREDICT\_US\_transport & micro & lipid\_contribution & profile\_rmse & 0.0066 & 0.0016 & 0.0118 & 2.7423 & 61.4000 & 0.0677 \\
PREDICT\_US\_transport & micro & lipid\_contribution & rmse & 0.0181 & 0.0098 & 0.0271 & 5.6815 & 65.4000 & 0.0185 \\
PREDICT\_US\_transport & micro & mac\_contribution & profile\_rmse & 0.0024 & 0.0003 & 0.0044 & 0.9898 & 59.4000 & 0.1021 \\
PREDICT\_US\_transport & micro & mac\_contribution & rmse & 0.0050 & 0.0020 & 0.0083 & 1.6457 & 57.2000 & 0.2745 \\
PREDICT\_US\_transport & micro & shrinkage\_vs\_raw\_mean & profile\_rmse & -0.0000 & -0.0000 & 0.0000 & -0.0000 & 35.8000 & 0.0195 \\
PREDICT\_US\_transport & micro & shrinkage\_vs\_raw\_mean & rmse & 0.0136 & 0.0077 & 0.0197 & 4.0410 & 63.2000 & 0.0285 \\
PREDICT\_US\_transport & micro & total\_gain\_context & profile\_rmse & 0.0080 & 0.0026 & 0.0135 & 3.3049 & 58.8000 & 0.1714 \\
PREDICT\_US\_transport & micro & total\_gain\_context & rmse & 0.0233 & 0.0084 & 0.0378 & 7.2179 & 61.0000 & 0.0810 \\
PREDICT\_US\_transport & population & channel\_increment\_PRIMARY & profile\_rmse & 0.0049 & 0.0024 & 0.0075 & 2.0097 & 60.0000 & 0.1413 \\
PREDICT\_US\_transport & population & channel\_increment\_PRIMARY & rmse & 0.0109 & 0.0048 & 0.0177 & 3.4919 & 62.8000 & 0.0185 \\
PREDICT\_US\_transport & population & full\_vs\_raw\_mean & profile\_rmse & 0.0049 & 0.0024 & 0.0075 & 2.0097 & 60.0000 & 0.1413 \\
PREDICT\_US\_transport & population & full\_vs\_raw\_mean & rmse & 0.0360 & 0.0181 & 0.0555 & 10.7026 & 59.6000 & 0.2115 \\
PREDICT\_US\_transport & population & level\_correction & profile\_rmse & -0.0000 & -0.0000 & 0.0000 & -0.0000 & 34.8000 & 0.0194 \\
PREDICT\_US\_transport & population & level\_correction & rmse & 0.0126 & -0.0019 & 0.0258 & 3.8870 & 61.0000 & 0.0565 \\
PREDICT\_US\_transport & population & lipid\_contribution & profile\_rmse & 0.0037 & 0.0013 & 0.0061 & 1.5146 & 60.8000 & 0.0821 \\
PREDICT\_US\_transport & population & lipid\_contribution & rmse & 0.0076 & 0.0027 & 0.0130 & 2.4597 & 60.8000 & 0.0569 \\
PREDICT\_US\_transport & population & mac\_contribution & profile\_rmse & 0.0011 & 0.0001 & 0.0022 & 0.4560 & 57.4000 & 0.2094 \\
PREDICT\_US\_transport & population & mac\_contribution & rmse & 0.0016 & 0.0001 & 0.0032 & 0.5165 & 54.6000 & 0.4541 \\
PREDICT\_US\_transport & population & shrinkage\_vs\_raw\_mean & profile\_rmse & -0.0000 & -0.0000 & 0.0000 & -0.0000 & 36.6000 & 0.0495 \\
PREDICT\_US\_transport & population & shrinkage\_vs\_raw\_mean & rmse & 0.0251 & 0.0130 & 0.0379 & 7.4733 & 58.8000 & 0.2080 \\
PREDICT\_US\_transport & population & total\_gain\_context & profile\_rmse & 0.0049 & 0.0024 & 0.0075 & 2.0097 & 60.0000 & 0.1413 \\
PREDICT\_US\_transport & population & total\_gain\_context & rmse & 0.0234 & 0.0131 & 0.0334 & 7.2453 & 67.0000 & 0.0039 \\
\bottomrule
\end{longtable}}

\clearpage

\section*{Supplementary Note 6: Prior-arm equivalence}
\addcontentsline{toc}{section}{Note 6: Prior-arm equivalence}

The equivalence margin was fixed in advance at 0.005, one third of the channel increment. In the
998-participant cohort, host conditioning beats an unconditional prior by $+0.00200$
(95\% CI, 0.00013--0.00380) and microbial conditioning by $+0.00181$ ($-0.00006$ to 0.00361).
Microbial against host conditioning gives $-0.00019$ ($-0.00058$ to 0.00020), which falls inside
the equivalence margin in all five draws. The bootstrap standard error implies a minimum
detectable advantage at 80\% power of 0.00056, which is 4\% of the channel increment. The
conclusion is equivalence with power, not an uninformative null.

{\footnotesize\setlength{\tabcolsep}{3pt}
\begin{longtable}{lllllllllll}
\toprule
\textbf{cohort} & \textbf{arm} & \textbf{reference arm} & \textbf{difference} & \textbf{lo} & \textbf{hi} & \textbf{relative pct} & \textbf{better} & \textbf{sign p} & \textbf{draws equivalent} & \textbf{margin} \\
\midrule
\endhead
PREDICT\_UK & host & population & 0.00200 & 0.00013 & 0.00380 & 0.66081 & 529.80000 & 0.09434 & 5 & 0.00500 \\
PREDICT\_UK & host\_micro & host & -0.00007 & -0.00032 & 0.00019 & -0.02226 & 482.80000 & 0.35866 & 5 & 0.00500 \\
PREDICT\_UK & host\_micro & micro & 0.00012 & -0.00016 & 0.00042 & 0.04130 & 503.40000 & 0.62677 & 5 & 0.00500 \\
PREDICT\_UK & micro & host & -0.00019 & -0.00058 & 0.00020 & -0.06359 & 494.60000 & 0.56817 & 5 & 0.00500 \\
PREDICT\_UK & micro & population & 0.00181 & -0.00006 & 0.00361 & 0.59767 & 530.40000 & 0.11130 & 5 & 0.00500 \\
PREDICT\_US\_transport & host & population & -0.00028 & -0.00593 & 0.00541 & -0.08991 & 49.40000 & 0.43986 & 0 & 0.00500 \\
PREDICT\_US\_transport & host\_micro & host & -0.00008 & -0.00054 & 0.00036 & -0.02685 & 48.60000 & 0.47924 & 5 & 0.00500 \\
PREDICT\_US\_transport & host\_micro & micro & -0.00027 & -0.00141 & 0.00081 & -0.08990 & 51.20000 & 0.55087 & 5 & 0.00500 \\
PREDICT\_US\_transport & micro & host & 0.00019 & -0.00086 & 0.00133 & 0.06295 & 46.40000 & 0.59011 & 5 & 0.00500 \\
PREDICT\_US\_transport & micro & population & -0.00009 & -0.00572 & 0.00547 & -0.02660 & 48.60000 & 0.43184 & 0 & 0.00500 \\
\bottomrule
\end{longtable}}

\clearpage

\section*{Supplementary Note 7: Classical baselines}
\addcontentsline{toc}{section}{Note 7: Classical baselines}

A linear mixed model with linear fixed effects, a random intercept and random channel slopes was
fitted on the training participants. For each held-out participant the best linear unbiased
predictor was computed from the identical calibration meals using the fitted variance components,
so the comparison differs only in the model, not in the information supplied.

In free living the mixed model reaches 2.532 without calibration and 2.156 with it, against 1.980
for the full model, a margin of 8.16\% (95\% CI, 0.061--0.295; 31 of 44 participants). In
standardized meals it reaches 0.369 and 0.357 against 0.299, a margin of 16.32\%
(95\% CI, 0.0518--0.0647; 737 of 998). Decomposing the margin, the nonlinear shared predictor
accounts for 10.76\% in the larger cohort, where it is estimated precisely; in the 44-participant
cohort the same comparison does not reach significance ($P=0.291$). The personal calibration gain
of the present method, 4.83\%, exceeds that of the mixed model's own random slopes, 3.27\%.

{\footnotesize\setlength{\tabcolsep}{3pt}
\begin{longtable}{llllll}
\toprule
\textbf{cohort} & \textbf{model} & \textbf{meals} & \textbf{people} & \textbf{rmse} & \textbf{profile rmse} \\
\midrule
\endhead
CGMacros & shared & 902 & 44 & 2.3639 & 1.8813 \\
CGMacros & raw\_personal\_mean & 902 & 44 & 2.1211 & 1.8813 \\
CGMacros & intercept\_only & 902 & 44 & 2.0420 & 1.8813 \\
CGMacros & full & 902 & 44 & 1.9802 & 1.8152 \\
CGMacros & lmm\_population & 902 & 44 & 2.5324 & 1.9956 \\
CGMacros & lmm\_blup & 902 & 44 & 2.1561 & 1.9395 \\
PREDICT\_UK & shared & 3987 & 998 & 0.3291 & 0.2482 \\
PREDICT\_UK & raw\_personal\_mean & 3987 & 998 & 0.3232 & 0.2482 \\
PREDICT\_UK & intercept\_only & 3987 & 998 & 0.3137 & 0.2482 \\
PREDICT\_UK & full & 3987 & 998 & 0.2985 & 0.2406 \\
PREDICT\_UK & lmm\_population & 3987 & 998 & 0.3688 & 0.2895 \\
PREDICT\_UK & lmm\_blup & 3987 & 998 & 0.3567 & 0.2841 \\
\bottomrule
\end{longtable}}

{\footnotesize\setlength{\tabcolsep}{3pt}
\begin{longtable}{lllllllllll}
\toprule
\textbf{cohort} & \textbf{contrast} & \textbf{method} & \textbf{reference} & \textbf{difference} & \textbf{diff lo} & \textbf{diff hi} & \textbf{relative pct} & \textbf{better} & \textbf{participants} & \textbf{sign p} \\
\midrule
\endhead
CGMacros & gutnut\_vs\_lmm\_blup & 1.9802 & 2.1561 & 0.1759 & 0.0611 & 0.2953 & 8.1595 & 31 & 44 & 0.0096 \\
CGMacros & gutnut\_vs\_lmm\_population & 1.9802 & 2.5324 & 0.5522 & 0.3421 & 0.7587 & 21.8064 & 35 & 44 & 0.0001 \\
CGMacros & lmm\_blup\_vs\_lmm\_population & 2.1561 & 2.5324 & 0.3763 & 0.1785 & 0.5918 & 14.8593 & 25 & 44 & 0.4514 \\
CGMacros & shared\_vs\_lmm\_population & 2.3639 & 2.5324 & 0.1685 & -0.0616 & 0.4265 & 6.6542 & 26 & 44 & 0.2912 \\
CGMacros & gutnut\_vs\_intercept\_only & 1.9802 & 2.0420 & 0.0618 & 0.0192 & 0.1036 & 3.0262 & 33 & 44 & 0.0013 \\
PREDICT\_UK & gutnut\_vs\_lmm\_blup & 0.2985 & 0.3567 & 0.0582 & 0.0518 & 0.0647 & 16.3150 & 737 & 998 & 0.0000 \\
PREDICT\_UK & gutnut\_vs\_lmm\_population & 0.2985 & 0.3688 & 0.0703 & 0.0633 & 0.0773 & 19.0501 & 761 & 998 & 0.0000 \\
PREDICT\_UK & lmm\_blup\_vs\_lmm\_population & 0.3567 & 0.3688 & 0.0120 & 0.0111 & 0.0130 & 3.2683 & 779 & 998 & 0.0000 \\
PREDICT\_UK & shared\_vs\_lmm\_population & 0.3291 & 0.3688 & 0.0397 & 0.0343 & 0.0451 & 10.7586 & 696 & 998 & 0.0000 \\
PREDICT\_UK & gutnut\_vs\_intercept\_only & 0.2985 & 0.3137 & 0.0152 & 0.0124 & 0.0179 & 4.8288 & 607 & 998 & 0.0000 \\
\bottomrule
\end{longtable}}

\clearpage

\section*{Supplementary Note 8: Identifiability simulation}
\addcontentsline{toc}{section}{Note 8: Identifiability simulation}

Personal coefficients were drawn from a known Gaussian population distribution with standard
deviations $(1.2,1.0,1.0)$, and meals were simulated with correlated fibre and fat exposures
mapped through the same saturating link. The grid crossed calibration length
$k\in\{1,3,5,10\}$, exposure correlation $\rho\in\{0,0.3,0.6,0.9\}$, residual scale
$\{0.5,1,2\}$ and prior displacement $\{0,1\}$ prior standard deviations, with 400 replicates per
cell.

With a correct prior, coverage at the 50\%, 80\% and 95\% levels remains near nominal at every
calibration length and correlation. Under a prior displaced by one prior standard deviation,
nominal 95\% coverage falls to 0.82--0.87. Coefficient recovery is poor throughout: at three meals
the recovery error for a channel slope is 0.935 against a prior standard deviation of 1.0, and the
posterior standard deviation shrinks by only 7.7\%. Transfer error to unseen compositions
nevertheless falls from 0.711 at one meal to 0.551 at three and 0.354 at ten. Three measured meals
therefore locate a participant within the population distribution without estimating their
coefficients, and personal adjustments must be read as model-based attributions.

{\footnotesize\setlength{\tabcolsep}{3pt}
\begin{longtable}{llllllllllll}
\toprule
\textbf{k} & \textbf{rho} & \textbf{sigma} & \textbf{prior bias} & \textbf{rmse mac} & \textbf{rmse lipid} & \textbf{transfer rmse} & \textbf{cov50 mac} & \textbf{cov80 mac} & \textbf{cov95 mac} & \textbf{shrink mac} & \textbf{median cond} \\
\midrule
\endhead
1 & 0.000 & 0.500 & 0.000 & 0.917 & 0.929 & 0.515 & 0.547 & 0.823 & 0.950 & 0.053 &  \\
3 & 0.000 & 0.500 & 0.000 & 0.819 & 0.878 & 0.351 & 0.510 & 0.802 & 0.963 & 0.143 & 26.733 \\
5 & 0.000 & 0.500 & 0.000 & 0.803 & 0.741 & 0.291 & 0.525 & 0.785 & 0.945 & 0.215 & 12.145 \\
10 & 0.000 & 0.500 & 0.000 & 0.662 & 0.618 & 0.221 & 0.477 & 0.795 & 0.945 & 0.344 & 9.024 \\
1 & 0.000 & 0.500 & 1.000 & 1.319 & 1.339 & 0.628 & 0.343 & 0.635 & 0.855 & 0.051 &  \\
3 & 0.000 & 0.500 & 1.000 & 1.169 & 1.166 & 0.427 & 0.338 & 0.637 & 0.845 & 0.150 & 23.671 \\
5 & 0.000 & 0.500 & 1.000 & 1.057 & 1.057 & 0.350 & 0.343 & 0.630 & 0.863 & 0.215 & 12.491 \\
10 & 0.000 & 0.500 & 1.000 & 0.839 & 0.856 & 0.261 & 0.440 & 0.715 & 0.895 & 0.340 & 9.002 \\
1 & 0.000 & 1.000 & 0.000 & 0.941 & 1.016 & 0.711 & 0.497 & 0.825 & 0.963 & 0.038 &  \\
3 & 0.000 & 1.000 & 0.000 & 0.935 & 0.913 & 0.550 & 0.522 & 0.797 & 0.943 & 0.077 & 25.065 \\
5 & 0.000 & 1.000 & 0.000 & 0.859 & 0.911 & 0.456 & 0.517 & 0.805 & 0.965 & 0.102 & 12.726 \\
10 & 0.000 & 1.000 & 0.000 & 0.853 & 0.874 & 0.354 & 0.470 & 0.777 & 0.955 & 0.164 & 9.056 \\
1 & 0.000 & 1.000 & 1.000 & 1.385 & 1.323 & 0.964 & 0.333 & 0.605 & 0.810 & 0.036 &  \\
3 & 0.000 & 1.000 & 1.000 & 1.319 & 1.279 & 0.605 & 0.350 & 0.595 & 0.818 & 0.073 & 24.752 \\
5 & 0.000 & 1.000 & 1.000 & 1.250 & 1.262 & 0.508 & 0.330 & 0.632 & 0.833 & 0.106 & 12.194 \\
10 & 0.000 & 1.000 & 1.000 & 1.141 & 1.095 & 0.410 & 0.357 & 0.647 & 0.840 & 0.164 & 8.948 \\
1 & 0.000 & 2.000 & 0.000 & 1.060 & 0.954 & 0.967 & 0.495 & 0.782 & 0.930 & 0.018 &  \\
3 & 0.000 & 2.000 & 0.000 & 0.937 & 1.011 & 0.773 & 0.485 & 0.805 & 0.963 & 0.039 & 25.164 \\
5 & 0.000 & 2.000 & 0.000 & 0.919 & 0.951 & 0.680 & 0.530 & 0.785 & 0.960 & 0.053 & 12.093 \\
10 & 0.000 & 2.000 & 0.000 & 0.866 & 0.958 & 0.532 & 0.517 & 0.843 & 0.963 & 0.078 & 9.059 \\
1 & 0.000 & 2.000 & 1.000 & 1.394 & 1.358 & 1.201 & 0.307 & 0.600 & 0.845 & 0.019 &  \\
3 & 0.000 & 2.000 & 1.000 & 1.386 & 1.308 & 0.946 & 0.333 & 0.618 & 0.828 & 0.039 & 26.794 \\
5 & 0.000 & 2.000 & 1.000 & 1.303 & 1.333 & 0.796 & 0.350 & 0.630 & 0.835 & 0.052 & 11.913 \\
10 & 0.000 & 2.000 & 1.000 & 1.253 & 1.318 & 0.650 & 0.340 & 0.623 & 0.860 & 0.077 & 9.029 \\
1 & 0.300 & 0.500 & 0.000 & 0.934 & 0.967 & 0.548 & 0.520 & 0.810 & 0.945 & 0.050 &  \\
3 & 0.300 & 0.500 & 0.000 & 0.915 & 0.870 & 0.359 & 0.463 & 0.765 & 0.932 & 0.140 & 27.122 \\
5 & 0.300 & 0.500 & 0.000 & 0.752 & 0.796 & 0.285 & 0.485 & 0.795 & 0.960 & 0.216 & 12.062 \\
10 & 0.300 & 0.500 & 0.000 & 0.590 & 0.646 & 0.217 & 0.550 & 0.848 & 0.965 & 0.343 & 8.819 \\
1 & 0.300 & 0.500 & 1.000 & 1.349 & 1.350 & 0.614 & 0.347 & 0.618 & 0.795 & 0.052 &  \\
3 & 0.300 & 0.500 & 1.000 & 1.211 & 1.227 & 0.446 & 0.333 & 0.625 & 0.825 & 0.136 & 26.615 \\
5 & 0.300 & 0.500 & 1.000 & 1.005 & 1.056 & 0.342 & 0.395 & 0.685 & 0.882 & 0.208 & 12.080 \\
10 & 0.300 & 0.500 & 1.000 & 0.811 & 0.820 & 0.257 & 0.395 & 0.713 & 0.877 & 0.339 & 9.071 \\
1 & 0.300 & 1.000 & 0.000 & 1.007 & 0.958 & 0.702 & 0.477 & 0.807 & 0.935 & 0.037 &  \\
3 & 0.300 & 1.000 & 0.000 & 0.894 & 0.888 & 0.494 & 0.525 & 0.833 & 0.945 & 0.077 & 25.405 \\
5 & 0.300 & 1.000 & 0.000 & 0.872 & 0.846 & 0.441 & 0.497 & 0.818 & 0.955 & 0.104 & 12.499 \\
10 & 0.300 & 1.000 & 0.000 & 0.870 & 0.849 & 0.342 & 0.485 & 0.780 & 0.930 & 0.167 & 8.957 \\
1 & 0.300 & 1.000 & 1.000 & 1.308 & 1.346 & 0.912 & 0.357 & 0.625 & 0.843 & 0.035 &  \\
3 & 0.300 & 1.000 & 1.000 & 1.322 & 1.326 & 0.608 & 0.367 & 0.632 & 0.835 & 0.073 & 24.926 \\
5 & 0.300 & 1.000 & 1.000 & 1.185 & 1.238 & 0.505 & 0.343 & 0.615 & 0.865 & 0.103 & 12.386 \\
10 & 0.300 & 1.000 & 1.000 & 1.083 & 1.178 & 0.418 & 0.378 & 0.670 & 0.860 & 0.165 & 8.788 \\
1 & 0.300 & 2.000 & 0.000 & 0.959 & 0.986 & 0.938 & 0.507 & 0.797 & 0.950 & 0.018 &  \\
3 & 0.300 & 2.000 & 0.000 & 0.965 & 0.955 & 0.766 & 0.537 & 0.800 & 0.940 & 0.041 & 25.401 \\
5 & 0.300 & 2.000 & 0.000 & 0.965 & 0.952 & 0.627 & 0.492 & 0.760 & 0.955 & 0.052 & 12.073 \\
10 & 0.300 & 2.000 & 0.000 & 0.937 & 0.962 & 0.558 & 0.482 & 0.767 & 0.948 & 0.077 & 9.004 \\
1 & 0.300 & 2.000 & 1.000 & 1.356 & 1.375 & 1.180 & 0.320 & 0.650 & 0.848 & 0.017 &  \\
3 & 0.300 & 2.000 & 1.000 & 1.332 & 1.330 & 0.953 & 0.340 & 0.610 & 0.828 & 0.039 & 24.801 \\
5 & 0.300 & 2.000 & 1.000 & 1.357 & 1.334 & 0.834 & 0.315 & 0.598 & 0.828 & 0.053 & 12.708 \\
10 & 0.300 & 2.000 & 1.000 & 1.305 & 1.206 & 0.622 & 0.357 & 0.603 & 0.825 & 0.078 & 8.714 \\
1 & 0.600 & 0.500 & 0.000 & 0.951 & 0.978 & 0.488 & 0.530 & 0.812 & 0.938 & 0.053 &  \\
3 & 0.600 & 0.500 & 0.000 & 0.861 & 0.842 & 0.335 & 0.512 & 0.795 & 0.950 & 0.143 & 26.011 \\
5 & 0.600 & 0.500 & 0.000 & 0.764 & 0.825 & 0.290 & 0.517 & 0.807 & 0.953 & 0.209 & 12.397 \\
10 & 0.600 & 0.500 & 0.000 & 0.683 & 0.690 & 0.224 & 0.527 & 0.785 & 0.927 & 0.337 & 8.818 \\
1 & 0.600 & 0.500 & 1.000 & 1.317 & 1.274 & 0.626 & 0.380 & 0.625 & 0.835 & 0.052 &  \\
3 & 0.600 & 0.500 & 1.000 & 1.162 & 1.173 & 0.424 & 0.347 & 0.657 & 0.845 & 0.136 & 26.049 \\
5 & 0.600 & 0.500 & 1.000 & 1.020 & 1.070 & 0.344 & 0.370 & 0.650 & 0.865 & 0.213 & 12.245 \\
10 & 0.600 & 0.500 & 1.000 & 0.819 & 0.853 & 0.262 & 0.390 & 0.688 & 0.905 & 0.334 & 9.057 \\
1 & 0.600 & 1.000 & 0.000 & 0.899 & 0.978 & 0.720 & 0.573 & 0.843 & 0.955 & 0.034 &  \\
3 & 0.600 & 1.000 & 0.000 & 0.894 & 0.871 & 0.519 & 0.502 & 0.802 & 0.958 & 0.075 & 24.885 \\
5 & 0.600 & 1.000 & 0.000 & 0.904 & 0.877 & 0.434 & 0.527 & 0.795 & 0.935 & 0.105 & 12.645 \\
10 & 0.600 & 1.000 & 0.000 & 0.824 & 0.877 & 0.344 & 0.527 & 0.802 & 0.963 & 0.160 & 9.053 \\
1 & 0.600 & 1.000 & 1.000 & 1.334 & 1.377 & 0.885 & 0.290 & 0.605 & 0.848 & 0.037 &  \\
3 & 0.600 & 1.000 & 1.000 & 1.312 & 1.299 & 0.631 & 0.360 & 0.598 & 0.815 & 0.073 & 24.788 \\
5 & 0.600 & 1.000 & 1.000 & 1.253 & 1.309 & 0.534 & 0.347 & 0.630 & 0.833 & 0.102 & 12.212 \\
10 & 0.600 & 1.000 & 1.000 & 1.125 & 1.094 & 0.409 & 0.335 & 0.645 & 0.858 & 0.159 & 9.022 \\
1 & 0.600 & 2.000 & 0.000 & 1.030 & 1.009 & 0.927 & 0.468 & 0.782 & 0.943 & 0.018 &  \\
3 & 0.600 & 2.000 & 0.000 & 0.927 & 0.939 & 0.791 & 0.555 & 0.805 & 0.948 & 0.039 & 26.611 \\
5 & 0.600 & 2.000 & 0.000 & 0.954 & 0.976 & 0.712 & 0.497 & 0.790 & 0.950 & 0.051 & 12.239 \\
10 & 0.600 & 2.000 & 0.000 & 0.974 & 0.967 & 0.545 & 0.490 & 0.775 & 0.943 & 0.077 & 9.181 \\
1 & 0.600 & 2.000 & 1.000 & 1.343 & 1.394 & 1.308 & 0.330 & 0.603 & 0.850 & 0.019 &  \\
3 & 0.600 & 2.000 & 1.000 & 1.253 & 1.410 & 0.901 & 0.323 & 0.655 & 0.870 & 0.039 & 25.663 \\
5 & 0.600 & 2.000 & 1.000 & 1.384 & 1.360 & 0.808 & 0.287 & 0.600 & 0.802 & 0.052 & 12.381 \\
10 & 0.600 & 2.000 & 1.000 & 1.266 & 1.282 & 0.617 & 0.340 & 0.620 & 0.820 & 0.076 & 8.917 \\
1 & 0.900 & 0.500 & 0.000 & 0.919 & 0.980 & 0.481 & 0.480 & 0.828 & 0.968 & 0.048 &  \\
3 & 0.900 & 0.500 & 0.000 & 0.873 & 0.875 & 0.345 & 0.525 & 0.815 & 0.925 & 0.135 & 36.620 \\
5 & 0.900 & 0.500 & 0.000 & 0.804 & 0.834 & 0.267 & 0.485 & 0.825 & 0.950 & 0.187 & 16.698 \\
10 & 0.900 & 0.500 & 0.000 & 0.705 & 0.675 & 0.213 & 0.497 & 0.810 & 0.958 & 0.294 & 12.096 \\
1 & 0.900 & 0.500 & 1.000 & 1.338 & 1.393 & 0.599 & 0.323 & 0.600 & 0.840 & 0.050 &  \\
3 & 0.900 & 0.500 & 1.000 & 1.173 & 1.179 & 0.383 & 0.335 & 0.632 & 0.858 & 0.135 & 38.810 \\
5 & 0.900 & 0.500 & 1.000 & 1.135 & 1.050 & 0.332 & 0.347 & 0.647 & 0.825 & 0.188 & 16.585 \\
10 & 0.900 & 0.500 & 1.000 & 0.860 & 0.925 & 0.243 & 0.390 & 0.682 & 0.895 & 0.296 & 11.882 \\
1 & 0.900 & 1.000 & 0.000 & 0.979 & 0.967 & 0.731 & 0.470 & 0.785 & 0.950 & 0.036 &  \\
3 & 0.900 & 1.000 & 0.000 & 0.943 & 0.938 & 0.534 & 0.505 & 0.797 & 0.948 & 0.070 & 36.401 \\
5 & 0.900 & 1.000 & 0.000 & 0.882 & 0.910 & 0.419 & 0.515 & 0.812 & 0.953 & 0.096 & 17.116 \\
10 & 0.900 & 1.000 & 0.000 & 0.831 & 0.859 & 0.331 & 0.477 & 0.812 & 0.955 & 0.146 & 12.102 \\
1 & 0.900 & 1.000 & 1.000 & 1.372 & 1.382 & 0.858 & 0.318 & 0.620 & 0.825 & 0.036 &  \\
3 & 0.900 & 1.000 & 1.000 & 1.353 & 1.246 & 0.597 & 0.333 & 0.585 & 0.825 & 0.070 & 33.489 \\
5 & 0.900 & 1.000 & 1.000 & 1.283 & 1.282 & 0.506 & 0.333 & 0.623 & 0.830 & 0.098 & 17.132 \\
10 & 0.900 & 1.000 & 1.000 & 1.183 & 1.173 & 0.398 & 0.355 & 0.625 & 0.855 & 0.146 & 12.558 \\
1 & 0.900 & 2.000 & 0.000 & 0.952 & 0.992 & 0.931 & 0.497 & 0.818 & 0.955 & 0.019 &  \\
3 & 0.900 & 2.000 & 0.000 & 1.002 & 0.972 & 0.821 & 0.530 & 0.777 & 0.932 & 0.039 & 33.530 \\
5 & 0.900 & 2.000 & 0.000 & 0.926 & 0.952 & 0.680 & 0.495 & 0.805 & 0.970 & 0.051 & 16.967 \\
10 & 0.900 & 2.000 & 0.000 & 0.949 & 0.957 & 0.575 & 0.510 & 0.782 & 0.945 & 0.072 & 12.152 \\
1 & 0.900 & 2.000 & 1.000 & 1.430 & 1.401 & 1.240 & 0.307 & 0.583 & 0.802 & 0.019 &  \\
3 & 0.900 & 2.000 & 1.000 & 1.395 & 1.414 & 0.946 & 0.318 & 0.585 & 0.800 & 0.038 & 35.455 \\
5 & 0.900 & 2.000 & 1.000 & 1.289 & 1.371 & 0.791 & 0.343 & 0.620 & 0.840 & 0.050 & 16.170 \\
10 & 0.900 & 2.000 & 1.000 & 1.329 & 1.267 & 0.671 & 0.335 & 0.598 & 0.828 & 0.073 & 12.352 \\
\bottomrule
\end{longtable}}

\clearpage

\section*{Supplementary Note 9: Uncertainty calibration, design conditioning and a second sensor}
\addcontentsline{toc}{section}{Note 9: Calibration and conditioning}

The residual noise variance is a single empirical-Bayes scalar shared by all participants, taking
values between 2.69 and 3.69 across outer folds. Three calibration meals and three coefficients
leave no residual degree of freedom from which a per-participant value could be estimated.

Interval honesty is assessed by posterior predictive coverage rather than by a calibration error
summary. Nominal 50\%, 80\% and 95\% intervals cover 0.659, 0.871 and 0.963 of observations, and
the standardized residuals have a standard deviation of 0.870 with a probability integral
transform mean of 0.482. Intervals are therefore conservative. The calibration slope is 0.884 and
the intercept 0.262, consistent with compression at high responses. Posterior standard deviations
fall by 36.4\%, 19.5\% and 17.5\% for the response level and the two channel slopes.

Three-meal calibration designs in free living are frequently ill conditioned. The median condition
number is 31.3 and 39\% of participants exceed 100. Prediction gain is nevertheless uncorrelated
with conditioning, which is consistent with a prior-dominated posterior.

Each meal carries a second continuous glucose sensor on 1{,}093 of 1{,}122 occasions, correlated
with the primary sensor at 0.865 with a median absolute difference of 0.679. This is a measurement
replication and not a second biological endpoint. Repeating the analysis on the second sensor gives
2.637 for the shared predictor, 2.335 for intercept-only calibration and 2.243 for the full model,
a channel increment of 0.0923 (95\% CI, 0.0313--0.1566; 27 of 44; $P=0.17$). The increment is
reproduced in magnitude while remaining an order of magnitude smaller than the disagreement
between the two sensors.

{\footnotesize\setlength{\tabcolsep}{3pt}
\begin{longtable}{lll}
\toprule
\textbf{block} & \textbf{quantity} & \textbf{value} \\
\midrule
\endhead
posterior\_predictive\_coverage & coverage\_50 & 0.6585 \\
posterior\_predictive\_coverage & coverage\_50\_participant\_balanced & 0.6588 \\
posterior\_predictive\_coverage & coverage\_80 & 0.8714 \\
posterior\_predictive\_coverage & coverage\_80\_participant\_balanced & 0.8741 \\
posterior\_predictive\_coverage & coverage\_95 & 0.9634 \\
posterior\_predictive\_coverage & coverage\_95\_participant\_balanced & 0.9643 \\
calibration & calibration\_slope & 0.8841 \\
calibration & calibration\_intercept & 0.2617 \\
calibration & pit\_mean & 0.4824 \\
calibration & pit\_sd & 0.2475 \\
calibration & standardised\_residual\_sd & 0.8702 \\
calibration\_design\_conditioning & median & 31.3147 \\
calibration\_design\_conditioning & q25 & 15.8129 \\
calibration\_design\_conditioning & q75 & 166075549731368608.0000 \\
calibration\_design\_conditioning & q90 & 166075549731368608.0000 \\
calibration\_design\_conditioning & max & 1980955931924112896.0000 \\
calibration\_design\_conditioning & share\_above\_100 & 0.3864 \\
calibration\_design\_conditioning & gain\_well\_conditioned\_median & 0.0134 \\
calibration\_design\_conditioning & gain\_ill\_conditioned\_median & 0.0365 \\
calibration\_design\_conditioning & spearman\_gain\_vs\_condition & -0.0026 \\
posterior\_contraction\_pct & intercept & 36.4500 \\
posterior\_contraction\_pct & mac & 19.5000 \\
posterior\_contraction\_pct & lipid & 17.5200 \\
second\_sensor\_dexcom & meals & 873.0000 \\
second\_sensor\_dexcom & participants & 44.0000 \\
second\_sensor\_dexcom & shared & 2.6365 \\
second\_sensor\_dexcom & intercept\_only & 2.3351 \\
second\_sensor\_dexcom & full & 2.2428 \\
second\_sensor\_dexcom\_channel\_increment & method & 2.2428 \\
second\_sensor\_dexcom\_channel\_increment & reference & 2.3351 \\
second\_sensor\_dexcom\_channel\_increment & difference & 0.0923 \\
second\_sensor\_dexcom\_channel\_increment & diff\_lo & 0.0313 \\
second\_sensor\_dexcom\_channel\_increment & diff\_hi & 0.1566 \\
second\_sensor\_dexcom\_channel\_increment & diff\_se & 0.0321 \\
second\_sensor\_dexcom\_channel\_increment & relative\_pct & 3.9531 \\
second\_sensor\_dexcom\_channel\_increment & rel\_lo & 1.4320 \\
second\_sensor\_dexcom\_channel\_increment & rel\_hi & 6.3387 \\
second\_sensor\_dexcom\_channel\_increment & participants & 44.0000 \\
second\_sensor\_dexcom\_channel\_increment & better & 27.0000 \\
second\_sensor\_dexcom\_channel\_increment & sign\_p & 0.1742 \\
shared\_residual\_variance & fold\_0\_sigma2 & 3.2493 \\
shared\_residual\_variance & fold\_1\_sigma2 & 3.3178 \\
shared\_residual\_variance & fold\_2\_sigma2 & 3.4412 \\
shared\_residual\_variance & fold\_3\_sigma2 & 2.6931 \\
shared\_residual\_variance & fold\_4\_sigma2 & 3.6928 \\
\bottomrule
\end{longtable}}

\clearpage

\section*{Supplementary Note 10: Cohort description, strata and responders}
\addcontentsline{toc}{section}{Note 10: Cohort description}

The free-living cohort comprises 14 participants with normal glucose tolerance, 16 with
prediabetes and 14 with type 2 diabetes. Median age is 51 years, median body mass index
30.2\,kg\,m$^{-2}$, median HbA1c 5.95\% and 28 of 44 participants are female. Of the meals, 6.6\%
have a positive incremental area of exactly zero. The standardized-meal cohorts have median ages
of 44 and 41 years and median body mass indices of 25.3 and 25.9\,kg\,m$^{-2}$, with 28.6\% of
meals at exactly zero.

The channel increment is positive in all three glycaemic strata, with median per-participant gains
of 0.028, 0.022 and 0.040 and 11 of 14, 12 of 16 and 10 of 14 participants improving.

{\footnotesize\setlength{\tabcolsep}{3pt}
\begin{longtable}{llllllllllllllll}
\toprule
\textbf{group} & \textbf{participants} & \textbf{meals} & \textbf{age median} & \textbf{age iqr} & \textbf{bmi median} & \textbf{bmi iqr} & \textbf{hba1c median} & \textbf{hba1c iqr} & \textbf{fasting glucose median} & \textbf{fasting glucose iqr} & \textbf{female n} & \textbf{iAUC median} & \textbf{iAUC iqr} & \textbf{iAUC zero pct} & \textbf{meals per person median} \\
\midrule
\endhead
all & 44 & 1122 & 51.00 & 43.75-58.25 & 30.16 & 26.94-35.93 & 5.95 & 5.50-6.90 & 110.50 & 100.00-142.50 & 28 & 1.94 & 0.667-3.899 & 6.60 & 25 \\
healthy & 14 & 365 & 40.00 & 28.50-48.50 & 26.46 & 23.42-30.72 & 5.35 & 5.12-5.47 & 97.00 & 91.50-101.50 & 9 & 1.16 & 0.442-2.252 & 6.03 & 25 \\
prediabetes & 16 & 410 & 54.00 & 45.75-58.25 & 32.76 & 28.83-36.22 & 5.95 & 5.80-6.12 & 110.50 & 104.00-125.25 & 11 & 2.04 & 0.779-3.718 & 4.63 & 26 \\
t2d & 14 & 347 & 52.50 & 51.00-59.00 & 30.62 & 28.24-34.35 & 7.10 & 6.93-7.35 & 149.00 & 142.50-157.00 & 8 & 3.22 & 1.474-5.909 & 9.51 & 25 \\
\bottomrule
\end{longtable}}

{\footnotesize\setlength{\tabcolsep}{3pt}
\begin{longtable}{llllllllll}
\toprule
\textbf{cohort} & \textbf{participants} & \textbf{families} & \textbf{meals} & \textbf{meal compositions} & \textbf{age median} & \textbf{bmi median} & \textbf{female n} & \textbf{target median} & \textbf{target zero pct} \\
\midrule
\endhead
PREDICT\_UK & 998 & 758 & 6981 & 7 & 44.00 & 25.30 & 651 & 0.54 & 28.63 \\
PREDICT\_US & 100 & 100 & 699 & 7 & 41.00 & 25.90 & 68 & 0.57 & 27.61 \\
\bottomrule
\end{longtable}}

\clearpage

\section*{Supplementary Note 11: Channel instantiation}
\addcontentsline{toc}{section}{Note 11: Channel instantiation}

The nutrient mask is a specification rather than a property of the method, so the choice of
exposures is an empirical question. Six instantiations were compared in the 998-participant cohort
with saturation constants set to the median non-zero exposure within each training fold.

Total fat is the best lipid coordinate of those tested. Full-model error is 0.3077 with a
twelve-item carbohydrate composite and total fat, and 0.3079 with fibre and total fat. Replacing
total fat by a fatty-acid fraction is worse in every case, reaching 0.3410 for saturated fat. The
channel increment itself becomes negative for the saturated and polyunsaturated instantiations,
meaning that personal calibration on those coordinates actively harms prediction. A composite of
twelve expert-assigned carbohydrate exposures performs equivalently to fibre alone.

The channel coordinate therefore needs to carry dose rather than compositional quality. This
conclusion is bounded by the composition space available: across only seven standardized meals,
fatty-acid fractions are nearly collinear with total fat and with meal identity.

{\footnotesize\setlength{\tabcolsep}{3pt}
\begin{longtable}{lllllll}
\toprule
\textbf{variant} & \textbf{model} & \textbf{meals} & \textbf{people} & \textbf{rmse} & \textbf{rmse sd} & \textbf{profile rmse} \\
\midrule
\endhead
fibre\_mufa & full & 3987.0000 & 998.0000 & 0.3140 & 0.0063 & 0.2447 \\
fibre\_mufa & intercept\_only & 3987.0000 & 998.0000 & 0.3198 & 0.0070 & 0.2459 \\
fibre\_mufa & intercept\_plus\_lipid & 3987.0000 & 998.0000 & 0.3099 & 0.0065 & 0.2405 \\
fibre\_mufa & intercept\_plus\_mac & 3987.0000 & 998.0000 & 0.3160 & 0.0066 & 0.2454 \\
fibre\_mufa & shared & 3987.0000 & 998.0000 & 0.3343 & 0.0046 & 0.2459 \\
fibre\_pufa & full & 3987.0000 & 998.0000 & 0.3276 & 0.0051 & 0.2578 \\
fibre\_pufa & intercept\_only & 3987.0000 & 998.0000 & 0.3169 & 0.0070 & 0.2459 \\
fibre\_pufa & intercept\_plus\_lipid & 3987.0000 & 998.0000 & 0.3222 & 0.0051 & 0.2534 \\
fibre\_pufa & intercept\_plus\_mac & 3987.0000 & 998.0000 & 0.3131 & 0.0067 & 0.2455 \\
fibre\_pufa & shared & 3987.0000 & 998.0000 & 0.3343 & 0.0046 & 0.2459 \\
fibre\_sfa & full & 3987.0000 & 998.0000 & 0.3409 & 0.0081 & 0.2691 \\
fibre\_sfa & intercept\_only & 3987.0000 & 998.0000 & 0.3157 & 0.0069 & 0.2459 \\
fibre\_sfa & intercept\_plus\_lipid & 3987.0000 & 998.0000 & 0.3360 & 0.0081 & 0.2643 \\
fibre\_sfa & intercept\_plus\_mac & 3987.0000 & 998.0000 & 0.3129 & 0.0067 & 0.2455 \\
fibre\_sfa & shared & 3987.0000 & 998.0000 & 0.3343 & 0.0046 & 0.2459 \\
fibre\_totalfat & full & 3987.0000 & 998.0000 & 0.3079 & 0.0065 & 0.2410 \\
fibre\_totalfat & intercept\_only & 3987.0000 & 998.0000 & 0.3189 & 0.0070 & 0.2459 \\
fibre\_totalfat & intercept\_plus\_lipid & 3987.0000 & 998.0000 & 0.3088 & 0.0065 & 0.2404 \\
fibre\_totalfat & intercept\_plus\_mac & 3987.0000 & 998.0000 & 0.3159 & 0.0067 & 0.2452 \\
fibre\_totalfat & shared & 3987.0000 & 998.0000 & 0.3343 & 0.0046 & 0.2459 \\
macmix\_sfa & full & 3987.0000 & 998.0000 & 0.3276 & 0.0075 & 0.2580 \\
macmix\_sfa & intercept\_only & 3987.0000 & 998.0000 & 0.3160 & 0.0069 & 0.2460 \\
macmix\_sfa & intercept\_plus\_lipid & 3987.0000 & 998.0000 & 0.3246 & 0.0075 & 0.2545 \\
macmix\_sfa & intercept\_plus\_mac & 3987.0000 & 998.0000 & 0.3127 & 0.0069 & 0.2453 \\
macmix\_sfa & shared & 3987.0000 & 998.0000 & 0.3343 & 0.0046 & 0.2460 \\
macmix\_totalfat & full & 3987.0000 & 998.0000 & 0.3077 & 0.0065 & 0.2408 \\
macmix\_totalfat & intercept\_only & 3987.0000 & 998.0000 & 0.3181 & 0.0070 & 0.2460 \\
macmix\_totalfat & intercept\_plus\_lipid & 3987.0000 & 998.0000 & 0.3093 & 0.0065 & 0.2412 \\
macmix\_totalfat & intercept\_plus\_mac & 3987.0000 & 998.0000 & 0.3166 & 0.0070 & 0.2457 \\
macmix\_totalfat & shared & 3987.0000 & 998.0000 & 0.3345 & 0.0046 & 0.2460 \\
\bottomrule
\end{longtable}}

{\footnotesize\setlength{\tabcolsep}{3pt}
\begin{longtable}{lllllllll}
\toprule
\textbf{variant} & \textbf{contrast} & \textbf{difference} & \textbf{lo} & \textbf{hi} & \textbf{relative pct} & \textbf{better} & \textbf{sign p} & \textbf{draws} \\
\midrule
\endhead
fibre\_mufa & channel\_increment & 0.0058 & 0.0022 & 0.0095 & 1.8092 & 510.0000 & 0.5287 & 3 \\
fibre\_mufa & lipid\_contribution & 0.0021 & -0.0012 & 0.0054 & 0.6534 & 502.6667 & 0.8421 & 3 \\
fibre\_mufa & mac\_contribution & -0.0041 & -0.0058 & -0.0024 & -1.3186 & 440.0000 & 0.0039 & 3 \\
fibre\_pufa & channel\_increment & -0.0107 & -0.0144 & -0.0069 & -3.3727 & 422.3333 & 0.0000 & 3 \\
fibre\_pufa & lipid\_contribution & -0.0145 & -0.0178 & -0.0112 & -4.6360 & 400.6667 & 0.0000 & 3 \\
fibre\_pufa & mac\_contribution & -0.0054 & -0.0069 & -0.0039 & -1.6669 & 425.6667 & 0.0009 & 3 \\
fibre\_sfa & channel\_increment & -0.0252 & -0.0291 & -0.0215 & -7.9756 & 336.0000 & 0.0000 & 3 \\
fibre\_sfa & lipid\_contribution & -0.0280 & -0.0319 & -0.0244 & -8.9647 & 310.3333 & 0.0000 & 3 \\
fibre\_sfa & mac\_contribution & -0.0049 & -0.0059 & -0.0040 & -1.4708 & 388.3333 & 0.0000 & 3 \\
fibre\_totalfat & channel\_increment & 0.0110 & 0.0084 & 0.0136 & 3.4412 & 569.0000 & 0.0001 & 3 \\
fibre\_totalfat & lipid\_contribution & 0.0079 & 0.0058 & 0.0102 & 2.5143 & 561.3333 & 0.0001 & 3 \\
fibre\_totalfat & mac\_contribution & 0.0009 & -0.0001 & 0.0019 & 0.2871 & 511.0000 & 0.1744 & 3 \\
macmix\_sfa & channel\_increment & -0.0115 & -0.0142 & -0.0089 & -3.6415 & 392.3333 & 0.0000 & 3 \\
macmix\_sfa & lipid\_contribution & -0.0148 & -0.0172 & -0.0126 & -4.7382 & 343.6667 & 0.0000 & 3 \\
macmix\_sfa & mac\_contribution & -0.0029 & -0.0038 & -0.0020 & -0.9005 & 404.0000 & 0.0000 & 3 \\
macmix\_totalfat & channel\_increment & 0.0104 & 0.0083 & 0.0127 & 3.2764 & 590.3333 & 0.0000 & 3 \\
macmix\_totalfat & lipid\_contribution & 0.0089 & 0.0070 & 0.0109 & 2.8120 & 596.6667 & 0.0000 & 3 \\
macmix\_totalfat & mac\_contribution & 0.0016 & 0.0011 & 0.0022 & 0.5257 & 560.0000 & 0.0029 & 3 \\
\bottomrule
\end{longtable}}

{\footnotesize\setlength{\tabcolsep}{3pt}
\begin{longtable}{llllll}
\toprule
\textbf{fold} & \textbf{variant} & \textbf{u channel1} & \textbf{u channel2} & \textbf{channel1} & \textbf{channel2} \\
\midrule
\endhead
0 & fibre\_totalfat & 1.850 & 22.200 & exposure\_\_Fiber, total dietary (g) & exposure\_\_Total Fat (g) \\
0 & fibre\_sfa & 1.850 & 3.496 & exposure\_\_Fiber, total dietary (g) & exposure\_\_Fatty acids, total saturated (g) \\
0 & fibre\_mufa & 1.850 & 9.500 & exposure\_\_Fiber, total dietary (g) & exposure\_\_Fatty acids, total monounsaturated (g) \\
0 & fibre\_pufa & 1.850 & 7.315 & exposure\_\_Fiber, total dietary (g) & exposure\_\_Fatty acids, total polyunsaturated (g) \\
0 & macmix\_totalfat & 8.716 & 22.200 & \_\_macmix\_\_ & exposure\_\_Total Fat (g) \\
0 & macmix\_sfa & 8.716 & 3.496 & \_\_macmix\_\_ & exposure\_\_Fatty acids, total saturated (g) \\
1 & fibre\_totalfat & 1.700 & 22.200 & exposure\_\_Fiber, total dietary (g) & exposure\_\_Total Fat (g) \\
1 & fibre\_sfa & 1.700 & 3.496 & exposure\_\_Fiber, total dietary (g) & exposure\_\_Fatty acids, total saturated (g) \\
1 & fibre\_mufa & 1.700 & 9.500 & exposure\_\_Fiber, total dietary (g) & exposure\_\_Fatty acids, total monounsaturated (g) \\
1 & fibre\_pufa & 1.700 & 7.315 & exposure\_\_Fiber, total dietary (g) & exposure\_\_Fatty acids, total polyunsaturated (g) \\
1 & macmix\_totalfat & 8.717 & 22.200 & \_\_macmix\_\_ & exposure\_\_Total Fat (g) \\
1 & macmix\_sfa & 8.717 & 3.496 & \_\_macmix\_\_ & exposure\_\_Fatty acids, total saturated (g) \\
2 & fibre\_totalfat & 1.850 & 22.200 & exposure\_\_Fiber, total dietary (g) & exposure\_\_Total Fat (g) \\
2 & fibre\_sfa & 1.850 & 3.496 & exposure\_\_Fiber, total dietary (g) & exposure\_\_Fatty acids, total saturated (g) \\
2 & fibre\_mufa & 1.850 & 9.500 & exposure\_\_Fiber, total dietary (g) & exposure\_\_Fatty acids, total monounsaturated (g) \\
2 & fibre\_pufa & 1.850 & 7.315 & exposure\_\_Fiber, total dietary (g) & exposure\_\_Fatty acids, total polyunsaturated (g) \\
2 & macmix\_totalfat & 8.716 & 22.200 & \_\_macmix\_\_ & exposure\_\_Total Fat (g) \\
2 & macmix\_sfa & 8.716 & 3.496 & \_\_macmix\_\_ & exposure\_\_Fatty acids, total saturated (g) \\
3 & fibre\_totalfat & 1.850 & 22.200 & exposure\_\_Fiber, total dietary (g) & exposure\_\_Total Fat (g) \\
3 & fibre\_sfa & 1.850 & 3.496 & exposure\_\_Fiber, total dietary (g) & exposure\_\_Fatty acids, total saturated (g) \\
3 & fibre\_mufa & 1.850 & 9.500 & exposure\_\_Fiber, total dietary (g) & exposure\_\_Fatty acids, total monounsaturated (g) \\
3 & fibre\_pufa & 1.850 & 7.315 & exposure\_\_Fiber, total dietary (g) & exposure\_\_Fatty acids, total polyunsaturated (g) \\
3 & macmix\_totalfat & 8.716 & 22.200 & \_\_macmix\_\_ & exposure\_\_Total Fat (g) \\
3 & macmix\_sfa & 8.716 & 3.496 & \_\_macmix\_\_ & exposure\_\_Fatty acids, total saturated (g) \\
4 & fibre\_totalfat & 2.000 & 22.200 & exposure\_\_Fiber, total dietary (g) & exposure\_\_Total Fat (g) \\
4 & fibre\_sfa & 2.000 & 3.496 & exposure\_\_Fiber, total dietary (g) & exposure\_\_Fatty acids, total saturated (g) \\
4 & fibre\_mufa & 2.000 & 9.500 & exposure\_\_Fiber, total dietary (g) & exposure\_\_Fatty acids, total monounsaturated (g) \\
4 & fibre\_pufa & 2.000 & 7.315 & exposure\_\_Fiber, total dietary (g) & exposure\_\_Fatty acids, total polyunsaturated (g) \\
4 & macmix\_totalfat & 8.716 & 22.200 & \_\_macmix\_\_ & exposure\_\_Total Fat (g) \\
4 & macmix\_sfa & 8.716 & 3.496 & \_\_macmix\_\_ & exposure\_\_Fatty acids, total saturated (g) \\
\bottomrule
\end{longtable}}

\clearpage

\section*{Supplementary Note 12: Evaluation protocols and the calibration curve}
\addcontentsline{toc}{section}{Note 12: Protocols and calibration curve}

Three protocols were compared. Under the fixed protocol every history length is scored on one
common query set. Under the deployment protocol the query set begins immediately after
calibration. Under the random protocol calibration meals are drawn at random rather than taken as
the first $k$.

The deployment protocol, which is the realistic one and scores more meals, gives a larger channel
increment than the fixed protocol (0.0736 against 0.0618 at three meals), so the fixed design was
not favourable. The random protocol, however, reduces the three-meal increment to 0.0343 with an
interval crossing zero ($P=0.55$); significance returns at five meals. Read together with the
condition numbers of Note 9, this identifies a requirement on the calibration set rather than a
failure of the method: measured meals must span the channel exposures. In standardized meals,
which span a wide composition range by design, the increment is strong under random calibration.

The calibration curve was extended to eleven meals so that the point at which it flattens is
observed. On the common query set, full-model error is 2.361 with the prior alone, 2.143 after one
meal, 2.053 after two, 1.936 after three, 1.919 after five, 1.891 after seven, 1.918 after nine
and 1.886 after eleven. Three meals capture approximately nine tenths of the improvement that
eleven meals reach. The channel increment is present at every history length, between 3.3\% and
5.5\% on the common query set.

{\footnotesize\setlength{\tabcolsep}{3pt}
\begin{longtable}{lllllll}
\toprule
\textbf{protocol} & \textbf{history k} & \textbf{model} & \textbf{meals} & \textbf{people} & \textbf{rmse} & \textbf{profile rmse} \\
\midrule
\endhead
deployment & 1 & full & 1078.0000 & 44.0000 & 2.1421 & 1.8045 \\
deployment & 1 & intercept\_only & 1078.0000 & 44.0000 & 2.2330 & 1.8633 \\
deployment & 1 & raw\_personal\_mean & 1078.0000 & 44.0000 & 3.0513 & 1.8633 \\
deployment & 1 & shared & 1078.0000 & 44.0000 & 2.3350 & 1.8633 \\
deployment & 3 & full & 990.0000 & 44.0000 & 1.9530 & 1.7984 \\
deployment & 3 & intercept\_only & 990.0000 & 44.0000 & 2.0266 & 1.8736 \\
deployment & 3 & raw\_personal\_mean & 990.0000 & 44.0000 & 2.1050 & 1.8736 \\
deployment & 3 & shared & 990.0000 & 44.0000 & 2.3445 & 1.8736 \\
deployment & 5 & full & 902.0000 & 44.0000 & 1.9625 & 1.8064 \\
deployment & 5 & intercept\_only & 902.0000 & 44.0000 & 2.0156 & 1.8813 \\
deployment & 5 & raw\_personal\_mean & 902.0000 & 44.0000 & 2.0537 & 1.8813 \\
deployment & 5 & shared & 902.0000 & 44.0000 & 2.3639 & 1.8813 \\
fixed & 1 & full & 902.0000 & 44.0000 & 2.1681 & 1.8266 \\
fixed & 1 & intercept\_only & 902.0000 & 44.0000 & 2.2616 & 1.8813 \\
fixed & 1 & raw\_personal\_mean & 902.0000 & 44.0000 & 3.0687 & 1.8813 \\
fixed & 1 & shared & 902.0000 & 44.0000 & 2.3639 & 1.8813 \\
fixed & 3 & full & 902.0000 & 44.0000 & 1.9802 & 1.8152 \\
fixed & 3 & intercept\_only & 902.0000 & 44.0000 & 2.0420 & 1.8813 \\
fixed & 3 & raw\_personal\_mean & 902.0000 & 44.0000 & 2.1211 & 1.8813 \\
fixed & 3 & shared & 902.0000 & 44.0000 & 2.3639 & 1.8813 \\
fixed & 5 & full & 902.0000 & 44.0000 & 1.9625 & 1.8064 \\
fixed & 5 & intercept\_only & 902.0000 & 44.0000 & 2.0156 & 1.8813 \\
fixed & 5 & raw\_personal\_mean & 902.0000 & 44.0000 & 2.0537 & 1.8813 \\
fixed & 5 & shared & 902.0000 & 44.0000 & 2.3639 & 1.8813 \\
random & 1 & full & 902.0000 & 44.0000 & 2.1961 & 1.8659 \\
random & 1 & intercept\_only & 902.0000 & 44.0000 & 2.2142 & 1.8810 \\
random & 1 & raw\_personal\_mean & 902.0000 & 44.0000 & 2.6868 & 1.8810 \\
random & 1 & shared & 902.0000 & 44.0000 & 2.3447 & 1.8810 \\
random & 3 & full & 902.0000 & 44.0000 & 2.0490 & 1.8458 \\
random & 3 & intercept\_only & 902.0000 & 44.0000 & 2.0833 & 1.8810 \\
random & 3 & raw\_personal\_mean & 902.0000 & 44.0000 & 2.2250 & 1.8810 \\
random & 3 & shared & 902.0000 & 44.0000 & 2.3447 & 1.8810 \\
random & 5 & full & 902.0000 & 44.0000 & 1.9936 & 1.8309 \\
random & 5 & intercept\_only & 902.0000 & 44.0000 & 2.0490 & 1.8810 \\
random & 5 & raw\_personal\_mean & 902.0000 & 44.0000 & 2.1189 & 1.8810 \\
random & 5 & shared & 902.0000 & 44.0000 & 2.3447 & 1.8810 \\
\bottomrule
\end{longtable}}

{\footnotesize\setlength{\tabcolsep}{3pt}
\begin{longtable}{lllllllll}
\toprule
\textbf{protocol} & \textbf{history k} & \textbf{difference} & \textbf{lo} & \textbf{hi} & \textbf{relative pct} & \textbf{better} & \textbf{sign p} & \textbf{draws} \\
\midrule
\endhead
deployment & 1 & 0.0909 & 0.0339 & 0.1483 & 4.0690 & 27.0000 & 0.1742 & 1 \\
deployment & 3 & 0.0736 & 0.0289 & 0.1160 & 3.6318 & 33.0000 & 0.0013 & 1 \\
deployment & 5 & 0.0531 & 0.0127 & 0.0964 & 2.6343 & 29.0000 & 0.0488 & 1 \\
fixed & 1 & 0.0935 & 0.0298 & 0.1579 & 4.1343 & 24.0000 & 0.6516 & 1 \\
fixed & 3 & 0.0618 & 0.0192 & 0.1036 & 3.0262 & 33.0000 & 0.0013 & 1 \\
fixed & 5 & 0.0531 & 0.0127 & 0.0964 & 2.6343 & 29.0000 & 0.0488 & 1 \\
random & 1 & 0.0181 & -0.0179 & 0.0590 & 0.8099 & 23.2000 & 0.7327 & 5 \\
random & 3 & 0.0343 & -0.0077 & 0.0765 & 1.6502 & 24.8000 & 0.5462 & 5 \\
random & 5 & 0.0554 & 0.0055 & 0.1059 & 2.6927 & 25.4000 & 0.4348 & 5 \\
\bottomrule
\end{longtable}}

{\footnotesize\setlength{\tabcolsep}{3pt}
\begin{longtable}{llllllll}
\toprule
\textbf{protocol} & \textbf{history k} & \textbf{profile rmse} & \textbf{profile rmse.1} & \textbf{profile rmse.2} & \textbf{rmse} & \textbf{rmse.1} & \textbf{rmse.2} \\
\midrule
\endhead
nan &  & full & intercept\_only & shared & full & intercept\_only & shared \\
common & 0.0000 & 1.8478 & 1.8526 & 1.8526 & 2.361 & 2.3699 & 2.3699 \\
common & 1.0000 & 1.7794 & 1.8526 & 1.8526 & 2.1427 & 2.2676 & 2.3699 \\
common & 2.0000 & 1.7778 & 1.8526 & 1.8526 & 2.0528 & 2.1232 & 2.3699 \\
common & 3.0000 & 1.757 & 1.8526 & 1.8526 & 1.9355 & 2.029 & 2.3699 \\
common & 5.0000 & 1.7399 & 1.8526 & 1.8526 & 1.9186 & 2.0026 & 2.3699 \\
common & 7.0000 & 1.7428 & 1.8526 & 1.8526 & 1.8909 & 1.9877 & 2.3699 \\
common & 9.0000 & 1.7447 & 1.8526 & 1.8526 & 1.9184 & 2.0081 & 2.3699 \\
common & 11.0000 & 1.755 & 1.8526 & 1.8526 & 1.8863 & 1.9774 & 2.3699 \\
deployment & 0.0000 & 1.8598 & 1.8633 & 1.8633 & 2.3296 & 2.335 & 2.335 \\
deployment & 1.0000 & 1.8045 & 1.8633 & 1.8633 & 2.1421 & 2.2329 & 2.335 \\
deployment & 2.0000 & 1.8105 & 1.8733 & 1.8733 & 2.0449 & 2.1026 & 2.3522 \\
deployment & 3.0000 & 1.7984 & 1.8736 & 1.8736 & 1.953 & 2.0266 & 2.3445 \\
deployment & 5.0000 & 1.8064 & 1.8813 & 1.8813 & 1.9625 & 2.0156 & 2.3639 \\
deployment & 7.0000 & 1.7807 & 1.8555 & 1.8555 & 1.9209 & 1.9869 & 2.3384 \\
deployment & 9.0000 & 1.7341 & 1.818 & 1.818 & 1.9386 & 1.9994 & 2.3574 \\
deployment & 11.0000 & 1.7408 & 1.834 & 1.834 & 1.8866 & 1.9699 & 2.3685 \\
\bottomrule
\end{longtable}}

{\footnotesize\setlength{\tabcolsep}{3pt}
\begin{longtable}{lllllllll}
\toprule
\textbf{protocol} & \textbf{history k} & \textbf{contrast} & \textbf{difference} & \textbf{diff lo} & \textbf{diff hi} & \textbf{relative pct} & \textbf{better} & \textbf{sign p} \\
\midrule
\endhead
common & 1 & channel\_increment & 0.1249 & 0.0397 & 0.2164 & 5.5078 & 32 & 0.0037 \\
common & 1 & total\_vs\_shared & 0.2272 & 0.0246 & 0.4029 & 9.5854 & 24 & 0.6516 \\
common & 2 & channel\_increment & 0.0704 & 0.0108 & 0.1268 & 3.3155 & 30 & 0.0226 \\
common & 2 & total\_vs\_shared & 0.3171 & 0.1410 & 0.4882 & 13.3805 & 25 & 0.4514 \\
common & 3 & channel\_increment & 0.0935 & 0.0282 & 0.1572 & 4.6091 & 31 & 0.0096 \\
common & 3 & total\_vs\_shared & 0.4344 & 0.2465 & 0.6049 & 18.3301 & 26 & 0.2912 \\
common & 5 & channel\_increment & 0.0839 & 0.0189 & 0.1558 & 4.1916 & 32 & 0.0037 \\
common & 5 & total\_vs\_shared & 0.4513 & 0.2825 & 0.6034 & 19.0429 & 30 & 0.0226 \\
common & 7 & channel\_increment & 0.0968 & 0.0198 & 0.1745 & 4.8693 & 31 & 0.0096 \\
common & 7 & total\_vs\_shared & 0.4790 & 0.2960 & 0.6415 & 20.2102 & 29 & 0.0488 \\
common & 9 & channel\_increment & 0.0897 & 0.0209 & 0.1577 & 4.4660 & 29 & 0.0488 \\
common & 9 & total\_vs\_shared & 0.4515 & 0.2642 & 0.6188 & 19.0508 & 27 & 0.1742 \\
common & 11 & channel\_increment & 0.0911 & 0.0239 & 0.1552 & 4.6080 & 26 & 0.2912 \\
common & 11 & total\_vs\_shared & 0.4836 & 0.2911 & 0.6565 & 20.4059 & 30 & 0.0226 \\
deployment & 1 & channel\_increment & 0.0909 & 0.0339 & 0.1483 & 4.0690 & 27 & 0.1742 \\
deployment & 1 & total\_vs\_shared & 0.1929 & -0.0223 & 0.3828 & 8.2626 & 29 & 0.0488 \\
deployment & 2 & channel\_increment & 0.0577 & 0.0137 & 0.1015 & 2.7430 & 30 & 0.0226 \\
deployment & 2 & total\_vs\_shared & 0.3073 & 0.1324 & 0.4802 & 13.0640 & 28 & 0.0961 \\
deployment & 3 & channel\_increment & 0.0736 & 0.0289 & 0.1160 & 3.6318 & 33 & 0.0013 \\
deployment & 3 & total\_vs\_shared & 0.3915 & 0.2190 & 0.5594 & 16.6970 & 28 & 0.0961 \\
deployment & 5 & channel\_increment & 0.0531 & 0.0127 & 0.0964 & 2.6343 & 29 & 0.0488 \\
deployment & 5 & total\_vs\_shared & 0.4014 & 0.2433 & 0.5478 & 16.9804 & 32 & 0.0037 \\
deployment & 7 & channel\_increment & 0.0660 & 0.0117 & 0.1171 & 3.3206 & 30 & 0.0226 \\
deployment & 7 & total\_vs\_shared & 0.4175 & 0.2389 & 0.5840 & 17.8523 & 28 & 0.0961 \\
deployment & 9 & channel\_increment & 0.0608 & 0.0080 & 0.1125 & 3.0403 & 30 & 0.0226 \\
deployment & 9 & total\_vs\_shared & 0.4188 & 0.2368 & 0.5845 & 17.7641 & 27 & 0.1742 \\
deployment & 11 & channel\_increment & 0.0833 & 0.0177 & 0.1476 & 4.2276 & 25 & 0.4514 \\
deployment & 11 & total\_vs\_shared & 0.4819 & 0.2879 & 0.6639 & 20.3445 & 28 & 0.0961 \\
\bottomrule
\end{longtable}}

\clearpage

\section*{Supplementary Note 13: Food scores against observed responses}
\addcontentsline{toc}{section}{Note 13: Food scores}

The personal response level shifts every food by the same amount and therefore cannot change any
within-participant ordering. This identity was verified numerically before the comparisons below,
so only the channel slopes can move an ordering.

Within-participant pairwise concordance with observed responses does not improve with calibration
in either cohort: $-0.0012$ (95\% CI, $-0.0103$ to 0.0075) in free living and $+0.0008$
($-0.0049$ to 0.0064) in standardized meals. Ties are frequent because each standardized-meal
participant contributes seven meals; 749 of 996 participants are exactly tied, and a sign test
that fails to exclude them is severely misleading.

Against a population nutritional-quality score the model orders observed responses 22.41
percentage points better (95\% CI, 21.03--23.77). The shared predictor alone already achieves
22.33 points of that margin, so composition modelling rather than personalization delivers it.

Of the total variance of the 44 by 1{,}226 score matrix, 69.6\% lies between participants, 8.4\%
between foods and 22.0\% in the person-by-food interaction. The mean absolute personal adjustment
is 2.94 score points. Personalization changes rankings substantially: the median rank shift is 94
places and a median of one food survives in a participant's top ten, while agreement with the
nutritional-quality reference falls slightly, from a median Spearman correlation of 0.433 to 0.397.

The library contains 1{,}226 items with FNDDS 2017-18 food codes expressed per 100\,kcal across 66
nutrient fields and 12 food groups, with a nutritional-quality reference for every item. Dairy and
beverages contribute only 4 and 11 items, so no group-level conclusion is drawn for them. Fibre is
zero for 12.0\% of items and fat for 0.9\%.

{\footnotesize\setlength{\tabcolsep}{3pt}
\begin{longtable}{llllllllllll}
\toprule
\textbf{cohort} & \textbf{comparison} & \textbf{mean a} & \textbf{mean b} & \textbf{difference} & \textbf{participants} & \textbf{better} & \textbf{worse} & \textbf{tied} & \textbf{sign p excluding ties} & \textbf{ci lo} & \textbf{ci hi} \\
\midrule
\endhead
CGMacros & full\_vs\_shared & 0.7058 & 0.7070 & -0.0012 & 44 & 20 & 21 & 3 & 1.0000 & -0.0103 & 0.0075 \\
PREDICT\_UK & full\_vs\_shared & 0.9259 & 0.9251 & 0.0008 & 996 & 125 & 122 & 749 & 0.8988 & -0.0049 & 0.0064 \\
PREDICT\_UK & full\_vs\_food\_compass & 0.9259 & 0.7018 & 0.2241 & 996 & 748 & 67 & 181 & 0.0000 & 0.2103 & 0.2377 \\
PREDICT\_UK & shared\_vs\_food\_compass & 0.9251 & 0.7018 & 0.2233 & 996 & 768 & 88 & 140 & 0.0000 & 0.2100 & 0.2370 \\
\bottomrule
\end{longtable}}

{\footnotesize\setlength{\tabcolsep}{3pt}
\begin{longtable}{lll}
\toprule
\textbf{block} & \textbf{quantity} & \textbf{value} \\
\midrule
\endhead
variance\_decomposition & cells & 53944 \\
variance\_decomposition & foods & 1226 \\
variance\_decomposition & participants & 44 \\
variance\_decomposition & food\_share\_pct & 8.369848874163965 \\
variance\_decomposition & participant\_share\_pct & 69.64700570341856 \\
variance\_decomposition & interaction\_share\_pct & 21.98314542241749 \\
variance\_decomposition & mean\_abs\_adjustment\_score\_points & 2.9367744134678775 \\
variance\_decomposition & sd\_adjustment\_score\_points & 6.752597576727614 \\
rank\_shift\_median & spearman & 0.8764288385791936 \\
rank\_shift\_median & median\_abs\_rank\_change & 94.25 \\
rank\_shift\_median & p90\_abs\_rank\_change & 287.0 \\
rank\_shift\_median & top10\_overlap & 1.0 \\
rank\_shift\_median & spearman\_vs\_fcs2\_shared & 0.4334703404451742 \\
rank\_shift\_median & spearman\_vs\_fcs2\_personal & 0.3973609874337672 \\
rank\_shift\_median & foods & 1226 \\
food\_library & foods & 1226 \\
food\_library & source & FNDDS 2017-18 food codes (food\_id prefix \\
food\_library & exposure\_basis & per 100 kcal \\
food\_library & food\_groups & 12 \\
food\_library & fcs2\_available & 1226 \\
food\_library & nutrient\_fields & 66 \\
food\_library & fibre\_zero\_pct & 11.99021207177814 \\
food\_library & fat\_zero\_pct & 0.897226753670473 \\
\bottomrule
\end{longtable}}

{\footnotesize\setlength{\tabcolsep}{3pt}
\begin{longtable}{ll}
\toprule
\textbf{food group} & \textbf{foods} \\
\midrule
\endhead
10000\_Beverages & 11 \\
1000\_Grains & 143 \\
2000\_Fruit & 57 \\
3000\_Vegetables & 374 \\
4000\_LegNut & 69 \\
5000\_MPE & 98 \\
5800\_Seafood & 83 \\
6000\_Dairy & 4 \\
7000\_Fats Oils & 23 \\
8000\_Mixed & 235 \\
8600\_SauceCondiment & 21 \\
9000\_SavorySweet & 108 \\
\bottomrule
\end{longtable}}

\clearpage

\section*{Supplementary Note 14: Cohort eligibility audit}
\addcontentsline{toc}{section}{Note 14: Cohort eligibility}

Candidate additional cohorts held locally were audited for eligibility as adaptive validation
data, which requires per-meal postprandial responses, microbial profiling and meal composition in
the same participants. None qualifies. Cross-sectional microbiome resources carry no meal-level
response. One resource carries per-meal responses but is recorded in the project provenance
registry as simulation material and is excluded from all evidence. The two cohorts used in this
work are therefore the complete set of eligible data.

{\footnotesize\setlength{\tabcolsep}{3pt}
\begin{longtable}{llll}
\toprule
\textbf{resource} & \textbf{per meal postprandial} & \textbf{cgm} & \textbf{meal composition} \\
\midrule
\endhead
CRA013939 (Zhang 2024 Adv Sci / CHNS) & False & False & False \\
predict\_multi L2 predict1\_real & True & True & False \\
external/american\_gut & False & False & False \\
gmrepo / cmd sqlite snapshots & False & False & False \\
verification/zoe2025 & False & False & False \\
\bottomrule
\end{longtable}}

\clearpage

\section*{Supplementary Figures}
\addcontentsline{toc}{section}{Supplementary Figures}

\noindent\textbf{Supplementary Fig. 1 $\vert$ Personal measurements reduce coefficient uncertainty.}\\
\includegraphics[width=\textwidth]{figure/Supplementary_Figure_1.pdf}

Prior and posterior coefficient standard deviations for the response level and the two channel
slopes in 44 participants after three measured meals. Lines pair estimates from the same
participant; black marks indicate medians. Median decreases are 36.4\%, 19.5\% and 17.5\%.
Standard deviations are conditional on the fitted model.

\clearpage
\noindent\textbf{Supplementary Fig. 2 $\vert$ Nutritional quality and response ordering.}\\
\includegraphics[width=\textwidth]{figure/Supplementary_Figure_2.pdf}

Participant-balanced error and within-participant ordering accuracy for an affine-calibrated
nutritional score, the shared predictor and a static channel model. Ordering uses native score
directions with half credit for ties.

\clearpage
\noindent\textbf{Supplementary Fig. 3 $\vert$ Error components and prior-information effects.}\\
\includegraphics[width=\textwidth]{figure/Supplementary_Figure_3.pdf}

Participant-balanced mean squared error partitioned into within-person variation and squared mean
error, and paired error reduction for the microbiome-informed prior relative to population and
host priors across calibration histories.

\clearpage
\noindent\textbf{Supplementary Fig. 4 $\vert$ Personal nutrient information beyond calibration compositions.}\\
\includegraphics[width=\textwidth]{figure/Supplementary_Figure_4.pdf}

Relative error reduction for all later meals, for compositions absent from personal calibration
and for channel exposure pairs absent from personal calibration, together with the distribution of
1{,}000 joint reassignments of the channel slopes between participants within the same outer fold.

\end{document}
